%% Discovery, diagnosis and repair of two physics defects in the MoCHII
%% Tier-1 collisional line-cooling fits, and a cooling-inventory comparison
%% against the other Monte Carlo photoionization codes.
%% Author: Kwang-Il Seon
%% Written for my_memo.cls.  Internal technical memo, not a paper.
\documentclass[english,a4paper]{my_memo}
\synctex=1
\usepackage{amsmath}
\usepackage{amssymb}
\usepackage{microtype}
\usepackage{booktabs}
\usepackage{array}
\hypersetup{hidelinks}

% my_memo.cls builds the trailing period of a section number into
% \thesection, so "Section~\ref{...}." would come out as "Section 3..".
% Move that period into the heading format instead, leaving \ref clean.
\makeatletter
\renewcommand{\thesection}{\@arabic\c@section}
\renewcommand{\thesubsection}{\thesection.\@arabic\c@subsection}
\renewcommand{\@seccntformat}[1]{\csname the#1\endcsname
  \@ifundefined{seccntdot@#1}{}{.}\quad}
\@namedef{seccntdot@section}{}
\makeatother

% float packing: these tables are small, let several share a page
\renewcommand{\topfraction}{0.92}
\renewcommand{\bottomfraction}{0.85}
\renewcommand{\textfraction}{0.07}
\renewcommand{\floatpagefraction}{0.72}
\setcounter{topnumber}{3}
\setcounter{totalnumber}{4}

\newcommand{\code}[1]{\texttt{#1}}
\newcommand{\MoCHII}{\textsc{MoCHII}}
\newcommand{\Cloudy}{\textsc{Cloudy}}
\newcommand{\Mocassin}{\textsc{Mocassin}}
\newcommand{\CMI}{\textsc{CMacIonize}}
\newcommand{\Torus}{\textsc{Torus}}
\newcommand{\HII}{H\,\textsc{ii}}
\newcommand{\HI}{H\,\textsc{i}}
\newcommand{\HeI}{He\,\textsc{i}}
\newcommand{\HeII}{He\,\textsc{ii}}
\newcommand{\HeIII}{He\,\textsc{iii}}
\newcommand{\OI}{O\,\textsc{i}}
\newcommand{\OII}{O\,\textsc{ii}}
\newcommand{\OIII}{O\,\textsc{iii}}
\newcommand{\NI}{N\,\textsc{i}}
\newcommand{\NII}{N\,\textsc{ii}}
\newcommand{\NIII}{N\,\textsc{iii}}
\newcommand{\CI}{C\,\textsc{i}}
\newcommand{\CII}{C\,\textsc{ii}}
\newcommand{\CIII}{C\,\textsc{iii}}
\newcommand{\CIV}{C\,\textsc{iv}}
\newcommand{\SII}{S\,\textsc{ii}}
\newcommand{\SIII}{S\,\textsc{iii}}
\newcommand{\SIV}{S\,\textsc{iv}}
\newcommand{\NeII}{Ne\,\textsc{ii}}
\newcommand{\NeIII}{Ne\,\textsc{iii}}
\newcommand{\Te}{T_{\rm e}}
\newcommand{\nel}{n_{\rm e}}
\newcommand{\OIV}{O\,\textsc{iv}}
\newcommand{\nH}{n_{\rm H}}
\newcommand{\ncrit}{n_{\rm crit}}

\begin{document}

\title{Two Physics Defects in the \MoCHII{} Line-Cooling Fits:\\
Discovery by a Fixed-State Cooling Budget, Repair,\\
and a Comparison with the Other Monte Carlo Codes}
\author{Kwang-Il Seon}
\date{\docmoddate}
\maketitle

\begin{abstract}
\MoCHII{} reproduced the ionization structure of the Lexington/Meudon
\HII{}-region benchmarks well but ran hotter than every published code:
$\langle T[N_{\rm p}N_{\rm e}]\rangle = 9159$~K at HII40 against a published
range of 7720--8199~K, and 7370~K at HII20 against 6402--6749~K.  Candidate
causes were excluded one at a time by direct experiment, which left the
cooling function itself.  A \emph{fixed-state} comparison --- evaluating
\MoCHII{}'s cooling expressions on \Cloudy{} c25.00's own converged zone
state, so that no feedback can hide a difference --- localized the problem
immediately: on an identical state the Tier-1 fits that drive the thermal
balance delivered $0.72$ of \Cloudy{}'s total cooling while the Tier-2
$n$-level solve that produces the line output delivered $0.98$.  The gap was
two defects in the offline fitting pipeline, both in \code{chianti\_cooling.py}
and its caller.  D1: the lower-level populations were distributed among the
ground-term fine-structure levels by Boltzmann factors, which at
$kT \gg \Delta E_{\rm fs}$ is the statistical-weight (high-density) limit and
suppressed the far-infrared fine-structure cooling ($\OIII{}$ $\times 1.64$,
$\NIII{}$ $\times 2.9$, $\CII{}$ $\times 2.5$ too small); volume-integrated
impact $23.9\%$ of the total cooling.  D2: the excitation energy was taken
from the \code{.scups} Burgess--Tully \emph{scaling} energy rather than the
observed \code{.elvlc} level difference, and entered both the Boltzmann factor
and the radiated photon ($\OII{}$: 30661 versus 26811~cm$^{-1}$, a factor
$0.583$ at 8000~K); impact $10.6\%$.  Missing channels were not the story
($\SIV{}$, the largest, is $1.2\%$).  Rebuilding all 28 fits from the
$\nel \to 0$ limit of statistical equilibrium with observed transition energies
moved HII40 to \textbf{7894~K, inside the published range}, and HII20 to
6376~K, $0.4\%$ below the published floor; $L({\rm H}\beta)$, $R_{\rm out}$
and the line ratios all moved into or next to the published spread at the same
time.  That zero-density limit still overcooled the low-critical-density lines
at nebular densities, leaving both benchmarks at the cool edge of the range;
suppressing the line cooling to the local $\nel$ with a $(T, \nel)$ table built
from the same $n$-level solve --- now the default,
\code{par\%cooling\_model = local\_ne} --- removes the far-infrared
overestimate, brings the fixed-state line total from $1.116\times$ to
$1.01\times$ \Cloudy{}'s, and superseded the two zero-density temperatures at
once by 8211 and 7025~K, matching \Cloudy{} c25.00 (8210, 6910~K); the
zero-density fits are retained as the \code{low\_density} option.  Corrections
to the diffuse field made after that carried the production figures to
\textbf{8169~K} at HII40 --- inside the published $7720$--$8199$~K range and
$0.50\%$ below \Cloudy{}'s 8210~K --- and \textbf{6925~K} at HII20, which sits
above the published $6402$--$6749$~K range but within $0.10\%$ of \Cloudy{}
c25.00's 6910~K, itself above that range: both current-generation codes now
run hotter than the 1995-era HII20 spread, and they agree with each other.
A comparison of the cooling
inventories of the other Monte Carlo photoionization codes is included; it
turned up defects on their side as well, and one omission common to all five
codes.
\end{abstract}

%==========================================================================
\section{Context: what was left after everything else was excluded}
\label{sec:context}

The Lexington/Meudon \HII{}-region benchmarks HII40 (40~kK blackbody) and
HII20 (20~kK blackbody) are the standard multi-code comparison for
photoionization codes; the published columns used here are those of Wood
et al.\ (2004, their Table~1), which agree entry by entry with the compilation
of Ercolano et al.\ (2003, their Tables~4 and 5).  Nine codes participate
(GF, HN, DP, TK, PH, RS, RR, BE and WME in the notation of that table), so
every ``published range'' quoted in this memo is the span of those nine
columns.  The \Cloudy{} c25.00 column, C25, is a run made for this work and is
deliberately \emph{not} counted in it: a 2025 code is not one of the original
participants, and keeping it separate is what makes it usable as an
independent current-generation reference.  \MoCHII{} matched the published
ionization structure but sat above every participating code in temperature:

\begin{center}
\begin{tabular}{lccc}
\toprule
 & \MoCHII{} (before) & published range & \Cloudy{} c25.00 \\
\midrule
HII40 $\langle T[N_{\rm p}N_{\rm e}]\rangle$ / K & 9159 & 7720--8199 & 8210 \\
HII20 $\langle T[N_{\rm p}N_{\rm e}]\rangle$ / K & 7370 & 6402--6749 & 6910 \\
\bottomrule
\end{tabular}
\end{center}

\noindent
that is $+14.1\%$ over the HII40 published median (8026~K) and $+11.6\%$ over
a current \Cloudy{} run of the same deck.  The following candidate causes were
each excluded by a direct experiment rather than by argument:

\begin{itemize}\itemsep2pt
\item \textbf{Charge exchange.}  A stage-mapping error was found and corrected
      (transition $i$ must use \code{chex}$(Z,i)$, not $(Z,i+1)$); the
      corrected rates changed HII40 $\Te$ by $-0.7\%$, not by $9\%$.
\item \textbf{Convergence of the reference.}  The stored \Mocassin{} baseline
      was underconverged; a full-quality re-run moved \emph{it} by $-230$~K
      and left the disagreement in place.
\item \textbf{Spatial resolution.}  A $25^3$ \Mocassin{} grid gives 8195~K
      against 8176~K on $13^3$.
\item \textbf{Recombination-coefficient vintage.}  Replacing the Badnell~2023
      metal RR+DR by \Mocassin{}'s Pequignot~1991 / Nussbaumer--Storey~1983
      set inside \MoCHII{} moved $\Te$ by $+0.2\%$ --- and in the wrong
      direction.
\item \textbf{\OIII{} collision strengths.}  The dominant coolant's
      $^3P \to {}^1D$ collision strength agrees between CHIANTI~v7 (the
      \Mocassin{} vintage) and v11 to $0.9\%$ (2.283 versus 2.262 at
      $10^4$~K), worth $<30$~K.
\item \textbf{Ionizing-band discretization.}  Aligning the band edges on the
      ionization thresholds fixed a real He and metal over-ionization but left
      $\Te$ unchanged.
\item \textbf{Transport treatment.}  Switching from case-B on-the-spot to the
      explicit Monte Carlo diffuse field \emph{raises} $\Te$ by $+2.6\%$
      (HII40) and $+1.6\%$ (HII20), so it cannot be the cause of an excess.
\item \textbf{Trace-species coupling.}  Adding the metal electrons to $\nel$
      and the metal photoheating to the balance moved HII40 from 9159 to
      9178~K, $+0.2\%$.
\item \textbf{Photoheating.}  The mean excess energy per photoionization is
      $\langle h\nu - 13.6\,{\rm eV}\rangle_{\rm H} = 3.9$~eV and
      $\langle h\nu - 24.6\,{\rm eV}\rangle_{\rm He} = 4.1$~eV, ordinary
      values for a mildly hardened 40~kK field, and the volumetric heating
      equals $\nel n_{\rm p}\alpha\langle h\nu - h\nu_{\rm th}\rangle$ with a
      standard $\alpha$.
\end{itemize}

A second apparent discrepancy turned out to be a consequence of the first
rather than an independent one.  $L({\rm H}\beta)$ was $12\%$ low at HII40.
Photon counting only holds for a radiation-bounded model: the temperature
excess lowers $\alpha_{\rm B}$ by $\times 0.897$, grows the Str\"omgren volume
by $\times 1.115$, and drives the ionization front past the prescribed outer
radius, so part of the ionizing luminosity leaves the model.  The predicted
leaked fraction, $10.3\%$, matched the measured one (recombinations
$3.788\times10^{49}$ against $Q_{\rm H} = 4.265\times10^{49}$~s$^{-1}$, an
$11.2\%$ escape), and HII20, which stays bounded, showed no such deficit.
So one open item remained, on the cooling side.

%==========================================================================
\section{The fixed-state cooling-budget method}
\label{sec:method}

Comparing two converged models confounds three things: the thermal balance,
the ionization structure, and the radiation field.  Each feeds back on the
others, and a cooling deficit hides itself by raising the temperature until
the books balance again.  The diagnosis therefore used a \emph{fixed-state}
comparison, in which no feedback is possible:

\begin{enumerate}\itemsep2pt
\item run \Cloudy{} c25.00 on the HII40 benchmark and save its zone-by-zone
      cooling inventory, heating inventory, overview state and ionization
      structure;
\item read the state $(\Te, \nel, \nH, x_{\rm ion})$ of every zone;
\item evaluate \MoCHII{}'s own cooling expressions on that state, channel by
      channel, and volume-integrate both sides on $4\pi r^2\,dr$.
\end{enumerate}

Any difference is then attributable to the atomic data and the cooling
formulas alone.  The \Cloudy{} run is the physics block of the distributed
test-suite deck \code{tsuite/auto/hii\_paris.in}, unmodified, with only output
commands appended (\code{tests/cloudy\_c25/hii40\_cool.in}); grains and cosmic
rays are absent from that deck, as its own log states.

Two evaluations of the \MoCHII{} side are carried through, because \MoCHII{}
computes metal line cooling twice, by two independent routes:
\begin{description}\itemsep2pt
\item[Tier-1] the fitted curves $\Lambda(T) = T^{-1/2}\sum_i A_i
      \exp(-T_i/T)$ read from \code{data/atomic/cooling\_<ion>.txt}.  These
      drive the thermal balance and are the object of this memo.
\item[Tier-2] the $n$-level statistical-equilibrium solve on the
      \code{nlevel\_<ion>.txt} models, used at output time for line
      luminosities.
\end{description}
Both are built from the same CHIANTI v11.0.2 files by the same offline
pipeline, so a disagreement between them on an identical state is internal
evidence, independent of \Cloudy{}.

\paragraph{Consistency checks performed before the comparison was believed.}
\begin{itemize}\itemsep2pt
\item Parsed ion fractions reproduce the \Cloudy{} file values exactly
      (\OIII{} 0.4570305, \SIV{} 0.0982921 in zone~1).
\item The sum of the ion cooling channels of an element reproduces that
      element's \code{.cooleach} column to $0.9886$--$0.9999$; the residual is
      the free-bound share that \code{.cooleach} also carries.
\item Free-free: the \code{.cool} ``FF'' entry equals the \code{.cooleach}
      ${\rm FF_H} + {\rm FF_M}$ sum to $1.0000$.
\item The \code{.ovr} \OIII{} column agrees with the \code{.ele\_o} $\rm O^{+2}$
      column to $4.4\times10^{-3}$.
\item \Cloudy{}'s own heating and cooling totals balance to $0.99987$, and its
      heating is $4.574$~eV per ionizing photon, split $90.8\%$ \HI{},
      $8.5\%$ \HeI{} --- the same physical picture as \MoCHII{}'s.
\item The Python Tier-2 solver used here reproduces \MoCHII{}'s own Fortran
      line luminosities on the same state to $0.1\%$ for all ten benchmark
      ions, so it is a faithful stand-in for the compiled code.
\end{itemize}

\paragraph{Result.}  Table~\ref{tab:budget} is the comparison.  The reading is
unambiguous: with the same state, the same database and the same code base,
the Tier-2 route reproduces \Cloudy{}'s total cooling to $2\%$ while the
Tier-1 fits that actually set the temperature deliver $72\%$ of it.  The
defect is therefore in the construction of the Tier-1 fits, not in the atomic
database, not in the missing channels, and not in the ionization structure.

\begin{table}[htbp]
\centering
\caption{HII40 volume-integrated cooling on the \Cloudy{} c25.00 state.
``Tier-1'' is \MoCHII{}'s fitted cooling as the thermal balance used it before
the fix; ``Tier-2'' replaces the metal line cooling by \MoCHII{}'s $n$-level
solve on the same state.  Both are ratios to the \Cloudy{} column.  The
free-bound rows are quoted case~A, which is why they exceed unity; \Cloudy{}
reports free-bound cooling net of escaping photons, so the matched bookkeeping
is the case-B line at the foot of the table.  Both \MoCHII{} columns are as
they stood at the time of the diagnosis, when \SIV{} had neither a cooling fit
nor an $n$-level model; it has both now (Section~\ref{sec:fix}).}
\label{tab:budget}
\begin{tabular}{lr@{\hspace{1.5em}}rcc}
\toprule
Channel & \Cloudy{} [erg\,s$^{-1}$] & \% of $L_{\rm cool}$ & Tier-1 & Tier-2 \\
\midrule
$[\OIII{}]$ lines            & $1.0020\times10^{38}$ & 32.06 & 0.587 & 1.024 \\
$[\OII{}]$ lines             & $4.8892\times10^{37}$ & 15.65 & 0.633 & 0.980 \\
$[\SIII{}]$ lines            & $4.7559\times10^{37}$ & 15.22 & 0.703 & 1.000 \\
H recombination (fb)         & $3.6809\times10^{37}$ & 11.78 & 1.635 & 1.635 \\
free-free (net)              & $2.3947\times10^{37}$ &  7.66 & 0.998 & 0.998 \\
$[\NII{}]$ lines             & $1.4818\times10^{37}$ &  4.74 & 0.922 & 0.990 \\
$[\SII{}]$ lines             & $6.7404\times10^{36}$ &  2.16 & 1.008 & 0.998 \\
$[\NeIII{}]$ lines           & $6.0443\times10^{36}$ &  1.93 & 0.933 & 1.299 \\
$[\NIII{}]$ lines            & $5.9069\times10^{36}$ &  1.89 & 0.433 & 1.019 \\
$[\CII{}]$ lines             & $5.5241\times10^{36}$ &  1.77 & 1.383 & 0.940 \\
$[\NeII{}]$ lines            & $4.1043\times10^{36}$ &  1.31 & 0.694 & 0.997 \\
$[\SIV{}]$ lines             & $3.8916\times10^{36}$ &  1.25 & 0.000 & 0.000 \\
\HeI{} recombination (fb)    & $3.4075\times10^{36}$ &  1.09 & 1.720 & 1.720 \\
\HI{} collisional lines      & $2.2490\times10^{36}$ &  0.72 & 0.989 & 0.989 \\
\CIII{}] lines               & $1.2741\times10^{36}$ &  0.41 & 0.822 & 0.990 \\
$[\OI{}]$ lines              & $3.1112\times10^{35}$ &  0.10 & 0.841 & 0.837 \\
\HeI{} collisional lines     & $1.9430\times10^{35}$ &  0.06 & 0.000 & 0.000 \\
metal recombination (fb)     & $1.6610\times10^{35}$ &  0.05 & 0.000 & 0.000 \\
\midrule
\textbf{Total (mapped)}      & $3.1244\times10^{38}$ & 99.98 & \textbf{0.819} & \textbf{1.077} \\
\textbf{Total, case-B recombination} & & & \textbf{0.723} & \textbf{0.981} \\
\bottomrule
\end{tabular}
\end{table}

Two further readings of Table~\ref{tab:budget} are worth recording.  First,
the ions whose Tier-1 ratio is farthest below unity ($\NIII{}$ 0.433,
$\OIII{}$ 0.587, $\OII{}$ 0.633, $\NeII{}$ 0.694, $\SIII{}$ 0.703) are exactly
the ions whose Tier-2 ratio is near unity, so the two tiers disagree with each
other and only one of them can be right.  Second, the missing channels are
small: $[\SIV{}]$ 10.5~$\mu$m, the largest one absent from \MoCHII{}
altogether, is $1.245\%$ of the budget, $[\OIV{}]$ is $0.007\%$, and every
other omission is below $0.2\%$.  Adding channels could not have closed a
$28\%$ gap.

%==========================================================================
\section{The two defects}
\label{sec:defects}

\subsection{D1: ground-term populations were a high-density model}

The fitting library built the cooling as
$\Lambda = \sum_{l,u} f_l\,q_{lu}(T)\,\Delta E_{lu}$ and set the lower-level
fractions $f_l$ by Boltzmann factors among the levels of the ground term:
\code{cooling\_effective} of
\code{tools/fitting/chianti\_cooling.py} called with the population model
\code{pop='ground\_term'}, from \code{tools/fitting/fit\_cooling.py}
(the call site is line~246 of that file, now reading
\code{cooling\_low\_density\_limit}).

For a ground term whose fine-structure splittings are far below $kT$ this is
wrong in a specific way.  At $T = 8000$~K, $kT = 5560$~cm$^{-1}$, while the
$\OIII{}$ $^3P$ levels lie at 0, 113 and 306~cm$^{-1}$: every Boltzmann factor
is within a few percent of unity, so the population collapses onto the
statistical weights, $g = 1, 3, 5$, and the lowest level receives $0.115$.
That is the LTE limit within the ground term, which is reached only for
$\nel \gg \ncrit$ of the fine structure.  The true population at the nebular
$\nel = 110$~cm$^{-3}$, where $\nel \ll \ncrit$, is the opposite limit: the
metastable fine-structure levels are drained radiatively faster than they are
collisionally populated, and the lowest level holds $\sim 0.75$.
Table~\ref{tab:pops} lists both for the benchmark ions.

The mechanism by which this loses cooling is the far-infrared fine-structure
lines.  The optical lines of $\OIII{}$ are excited out of all three $^3P$
levels, so their total is nearly insensitive to how the $^3P$ population is
distributed (the fractions sum to one either way).  The $[\OIII{}]$ 52 and
88~$\mu$m lines, however, are excitations \emph{within} the term, out of the
lowest levels into the higher ones; emptying the lowest level empties them.
At the benchmark state the $n$-level solve puts $46.7\%$ of the $\OIII{}$
cooling in 52 + 88~$\mu$m and $53.3\%$ in the optical lines --- and the
defective Tier-1 fit returned $0.594$ of the Tier-2 value, that is, almost
exactly the optical share alone.

Volume-integrated over the HII40 model, D1 accounts for $23.9\%$ of the total
cooling.

\begin{table}[htbp]
\centering
\caption{D1: the ground-term population model.  ``Boltzmann'' is the fraction
in the lowest level under the population model that was used; ``true'' is the
fraction from the $n$-level statistical equilibrium at $\nel = 110$~cm$^{-3}$.
The last column is the old Tier-1 fit divided by the Tier-2 solve at the same
state; it exceeds unity for $\CII{}$ because the $[\CII{}]$ 158~$\mu$m line is
already collisionally saturated at this density (Section~\ref{sec:limit}).
$T = 8000$~K.}
\label{tab:pops}
\begin{tabular}{llccc}
\toprule
Ion & ground-term levels $(g, E\,[{\rm cm^{-1}}])$ & Boltzmann $f_1$ & true $f_1$ & Tier-1/Tier-2 \\
\midrule
$\OIII{}$  & (1, 0) (3, 113) (5, 306)      & 0.115 & 0.752 & 0.594 \\
$\SIII{}$  & (1, 0) (3, 298) (5, 833)      & 0.123 & 0.907 & 0.700 \\
$\NII{}$   & (1, 0) (3, 49) (5, 131)       & 0.113 & 0.280 & 0.905 \\
$\NIII{}$  & (2, 0) (4, 174)               & 0.340 & 0.885 & 0.425 \\
$\CII{}$   & (2, 0) (4, 63)                & 0.336 & 0.414 & 1.933 \\
$\NeII{}$  & (4, 0) (2, 780)               & 0.697 & 1.000 & 0.698 \\
$\NeIII{}$ & (5, 0) (3, 643) (1, 920)      & 0.587 & 0.999 & 0.731 \\
$\OII{}$   & (4, 0) (4, 26831) (6, 26811)  & 0.980 & 1.000 & 0.596 \\
$\SII{}$   & (4, 0) (4, 14853) (6, 14885)  & 0.853 & 0.996 & 1.017 \\
$\OI{}$    & (5, 0) (3, 158) (1, 227)      & 0.563 & 0.999 & 1.002 \\
\bottomrule
\end{tabular}
\end{table}

\subsection{D2: the Burgess--Tully scaling energy was used as a physical
energy}

The excitation energy of every transition was taken from the \code{.scups}
header field, \code{de\_erg = tr["de"] * RY\_ERG}.  That field is the
Burgess \& Tully \emph{scaling} energy: a theoretical value whose only role is
to build the scaled temperature on which the collision strength is tabulated.
It is not the physical transition energy, and it entered the calculation
twice --- in the Boltzmann factor $\exp(-\Delta E/kT)$ of the excitation rate
and as the energy of the radiated photon.

The physical requirement is that both be the observed \code{.elvlc} level
difference.  Detailed balance ties the collisional excitation rate to the same
level gap that appears in the statistical-equilibrium matrix; if the cooling
sum and the $n$-level solve use different gaps for the same transition, they
are describing different atoms.  For $[\OII{}]$ 3726/3729 the \code{.scups}
value is 30661~cm$^{-1}$ against an observed 26811~cm$^{-1}$: the ratio of
Boltzmann factors alone is $\exp(-3850\,{\rm cm^{-1}}\,hc/kT) = 0.50$ at
8000~K, and with the smaller radiated energy the net factor is $0.583$.

Volume-integrated over HII40, D2 accounts for $10.6\%$ of the total cooling.

\begin{table}[htbp]
\centering
\caption{Isolating the two defects.  $\Lambda$ at $T = 8000$~K
[erg\,cm$^3$\,s$^{-1}$] per $\nel n_{\rm ion}$, with each defect repaired
separately and then together.  ``both'' must reproduce the $\nel \to 0$ limit
of the $n$-level solve, and does.  The last column is the $n$-level solve at
the nebular $\nel = 110$~cm$^{-3}$, which additionally carries the density
suppression the fits do not (Section~\ref{sec:limit}).}
\label{tab:isolate}
\begin{tabular}{lcccc@{\hspace{2em}}cc}
\toprule
 & \multicolumn{4}{c}{$\Lambda(8000\,{\rm K})$} & \multicolumn{2}{c}{$n$-level solve} \\
\cmidrule(r){2-5}\cmidrule(l){6-7}
Ion & as fitted & D2 fixed & D1 fixed & both & $\nel \to 0$ & $\nel = 110$ \\
\midrule
$\OIII{}$  & 2.691e$-$21 & 3.328e$-$21 & 4.401e$-$21 & 4.980e$-$21 & 4.970e$-$21 & 4.564e$-$21 \\
$\OII{}$   & 8.465e$-$22 & 1.452e$-$21 & 8.465e$-$22 & 1.452e$-$21 & 1.452e$-$21 & 1.421e$-$21 \\
$\SIII{}$  & 3.179e$-$20 & 3.179e$-$20 & 4.691e$-$20 & 4.690e$-$20 & 4.691e$-$20 & 4.543e$-$20 \\
$\NII{}$   & 4.988e$-$21 & 5.655e$-$21 & 5.470e$-$21 & 6.155e$-$21 & 6.108e$-$21 & 5.512e$-$21 \\
$\NIII{}$  & 7.412e$-$22 & 7.251e$-$22 & 2.140e$-$21 & 2.109e$-$21 & 2.105e$-$21 & 1.742e$-$21 \\
$\NeII{}$  & 7.019e$-$22 & 7.022e$-$22 & 1.007e$-$21 & 1.007e$-$21 & 1.007e$-$21 & 1.007e$-$21 \\
$\NeIII{}$ & 2.143e$-$21 & 2.143e$-$21 & 2.932e$-$21 & 2.933e$-$21 & 2.932e$-$21 & 2.929e$-$21 \\
$\CII{}$   & 5.892e$-$22 & 5.891e$-$22 & 1.450e$-$21 & 1.450e$-$21 & 1.366e$-$21 & 3.048e$-$22 \\
$\SII{}$   & 3.458e$-$20 & 3.470e$-$20 & 3.458e$-$20 & 3.470e$-$20 & 3.468e$-$20 & 3.396e$-$20 \\
$\OI{}$    & 4.702e$-$22 & 4.702e$-$22 & 4.693e$-$22 & 4.693e$-$22 & 4.691e$-$22 & 4.691e$-$22 \\
$\CIII{}$  & 4.724e$-$23 & 5.887e$-$23 & 4.724e$-$23 & 5.887e$-$23 & 5.845e$-$23 & 5.844e$-$23 \\
\bottomrule
\end{tabular}
\end{table}

\subsection{Temperature closure}

The two defects together are worth $34.1\%$ of the HII40 cooling.  The net
shortfall of the fits against the $n$-level solve is smaller, $25.8\%$,
because the low-density limit gives back $-8.3\%$ through the density
suppression it does not carry.

Whether that deficit explains the temperature excess can be checked without
running the code, because at fixed state the cooling is a smooth power of $T$.
The logarithmic slope of the total \MoCHII{} cooling at the \Cloudy{} state is
$d\ln\Lambda/d\ln T = 1.748$.  Scaling $\Te$ by a factor $f$ until \MoCHII{}'s
cooling equals \Cloudy{}'s --- \Cloudy{} balances its own heating at $f = 1$,
and the two codes' photoheating agrees --- gives
\[
  f = 1.111 \;\longrightarrow\; \Te = 9118~{\rm K},
\]
against the 9159~K \MoCHII{} actually converged to and the 8210~K of
\Cloudy{}.  Repeating with case-B recombination bookkeeping gives 9544~K, so
the estimate brackets the observed value.  In other words the cooling defect
accounts for essentially the entire $+11.6\%$ excess.

A second, independent estimate uses \MoCHII{}'s own converged state: there,
the Tier-2 cooling is $1.270\times$ the Tier-1 cooling at the same
temperature, and equals it at $f = 0.856$, that is $\langle \Te\rangle =
7836$~K in place of 9159~K.  The two estimates agree that repairing the fits
moves the benchmark from above the published range to inside it.

%==========================================================================
\section{The fix}
\label{sec:fix}

\subsection{The right population model}

\code{cooling\_low\_density\_limit()}
(\code{tools/fitting/chianti\_cooling.py}, line~232) replaces the ground-term
Boltzmann model with the $\nel \to 0$ limit of the $n$-level statistical
equilibrium,
\[
  \Lambda(T) \;=\; \sum_u q_{1u}(T)\,\Delta E_{1u},
\]
the sum running over every \code{.scups} transition out of the lowest level.
The physical argument is the one the defect got backwards: at
$\nel \ll \ncrit$ every ion sits in its lowest level, so collisional excitation
happens out of that level only, and each excitation is followed by a radiative
cascade that ultimately radiates the full $\Delta E_{1u}$ whatever route it
takes down.  The total cooling is therefore the excitation power out of the
ground level, which is exactly what the $n$-level solve returns as
$\nel \to 0$.  That identity is the strongest verification available and is
checked directly (Section~\ref{sec:fitver}).

\subsection{The right energies}

A new function, \code{transition\_energy\_erg()} at line~213 of the same file,
returns the observed \code{.elvlc} level difference.
\code{cooling\_effective} now calls it at line~333, in place of the old
\code{de\_erg = tr["de"] * RY\_ERG}.  Transitions whose observed
energies do not order the two levels return zero and are skipped rather than
silently contributing a negative energy.  The \code{.scups} \code{de} field is
retained for its only legitimate use, building the scaled temperature.

\subsection{A fit ladder derived from the atomic structure}

The repair also removed a piece of hand tuning that had been hiding the
defect.  In the low-density limit the cooling is
\[
  \Lambda(T) \;=\; T^{-1/2}\sum_u
      \Big[\tfrac{C}{g_1}\,\Upsilon_{1u}(T)\,\Delta E_{1u}\Big]\,
      e^{-T_u/T}, \qquad T_u = \Delta E_{1u}/k,
\]
which \emph{is} the fitted form $\Lambda = T^{-1/2}\sum_i A_i e^{-T_i/T}$
apart from the slow temperature dependence of $\Upsilon$.  The fitted $T_i$
are therefore excitation temperatures, not free parameters, and
\code{ground\_excitation\_ladder()} (\code{fit\_cooling.py}, line~115) now
builds the starting ladder by clustering the transitions in $\log T_u$ with
the share of the cooling each carries as the weight.  The hand-tuned guess
lists that the previous version carried for each ion are gone, which is why
the same script can now be pointed at a new ion without tuning.

\subsection{Verification of the rebuilt fits}
\label{sec:fitver}

Three independent checks were run on all 28 rebuilt files
(Table~\ref{tab:newfits}):

\begin{enumerate}\itemsep2pt
\item \textbf{Fit against its own target.}  The fitted curve against the
      directly evaluated $\nel \to 0$ statistical-equilibrium cooling over
      $10^3$--$10^5$~K.  The H/C/N/O/Ne/S set is within $0.7\%$ over the whole
      range except $\CIII{}$ ($5.5\%$ at the cold end, $0.93\%$ at $10^4$~K)
      and $\SII{}$ ($2.2\%$); the depleted Si, Cl and Ca ions reach $4$--$28\%$
      at the ends of the range, which is accepted because at abundances
      $\sim 3\times10^{-7}$ their share of the thermal balance is $\sim
      10^{-4}$ and their diagnostics run through the exact Tier-2 solve.
\item \textbf{Fit against a different code path.}  The fitted curve against
      the Tier-2 $n$-level solve at $\nel = 1$~cm$^{-3}$ --- a different data
      file and a different solver.  Agreement is $0.1$--$0.2\%$ for the
      benchmark ions.  The outliers are informative rather than alarming:
      Fe~\textsc{ii} $0.873$, Fe~\textsc{iii} $0.842$ and Ca~\textsc{v}
      $1.045$ differ because their $n$-level models are level-truncated, while
      the Tier-1 fit sums the full \code{.scups} list.
\item \textbf{New against old.}  The ion-by-ion multiplier at 5, 8 and
      10~kK, which is the quantitative statement of how much cooling the
      defects were losing.
\end{enumerate}

\begin{table}[htbp]
\centering
\caption{The rebuilt Tier-1 fits.  ``max err'' and ``err at $10^4$~K'' are
against the directly evaluated $\nel\to0$ statistical-equilibrium cooling over
$10^3$--$10^5$~K; ``new/old'' is the rebuilt fit divided by the defective one;
the last column is the independent Tier-2 solve at $\nel = 1$~cm$^{-3}$ divided
by the new fit.  $\SIV{}$ is new, so it has no old file.  The Ar, Mg, Fe, Si,
Cl and Ca ions are omitted here; their multipliers run 0.66--2.83 and their
default abundances are zero.}
\label{tab:newfits}
\begin{tabular}{lcc@{\hspace{2em}}ccc@{\hspace{2em}}c}
\toprule
 & \multicolumn{2}{c}{fit quality} & \multicolumn{3}{c}{new/old} & Tier-2$(\nel{=}1)$ \\
\cmidrule(r){2-3}\cmidrule(r){4-6}\cmidrule(l){7-7}
Ion & max & at $10^4$~K & 5~kK & 8~kK & 10~kK & / new \\
\midrule
\HI{}      & 0.52\% & 0.20\% & 1.012 & 1.008 & 1.006 & 1.000 \\
C\,\textsc{i}   & 0.27\% & 0.02\% & 1.189 & 1.049 & 1.017 & --- \\
$\CII{}$   & 0.35\% & 0.12\% & 2.956 & 2.460 & 1.742 & 0.943 \\
$\CIII{}$  & 5.48\% & 0.93\% & 1.442 & 1.246 & 1.191 & 0.971 \\
$\NI{}$    & 0.58\% & 0.14\% & 1.161 & 1.092 & 1.069 & --- \\
$\NII{}$   & 0.66\% & 0.31\% & 1.557 & 1.233 & 1.163 & 0.993 \\
$\NIII{}$  & 0.70\% & 0.23\% & 2.877 & 2.845 & 2.564 & 0.998 \\
$\OI{}$    & 0.20\% & 0.08\% & 1.035 & 0.998 & 0.997 & 1.000 \\
$\OII{}$   & 0.27\% & 0.07\% & 2.616 & 1.716 & 1.492 & 1.000 \\
$\OIII{}$  & 0.32\% & 0.07\% & 2.731 & 1.836 & 1.501 & 0.999 \\
$\NeII{}$  & 0.40\% & 0.12\% & 1.397 & 1.432 & 1.446 & 1.001 \\
$\NeIII{}$ & 0.34\% & 0.11\% & 1.501 & 1.369 & 1.247 & 1.001 \\
$\SII{}$   & 2.23\% & 0.96\% & 1.005 & 1.003 & 1.003 & 1.001 \\
$\SIII{}$  & 0.23\% & 0.01\% & 1.944 & 1.474 & 1.349 & 1.001 \\
$\SIV{}$   & 0.61\% & 0.16\% & \multicolumn{3}{c}{(new ion)}   & 1.002 \\
\bottomrule
\end{tabular}
\end{table}

\subsection{\SIV{} added}

$[\SIV{}]$ 10.51~$\mu$m was the one channel of any size that \MoCHII{} did not
carry at all.  \SIV{} was already tracked in the ionization cascade, so adding
it was the intended data operation: \code{cooling\_s\_4.txt} (three
exponential terms, max fit error $0.61\%$) and \code{nlevel\_s\_4.txt} from
the same pipeline.  In the rebuilt HII40 model $[\SIV{}]$ 10.51~$\mu$m comes
out at $0.1868$ of ${\rm H}\beta$ against $0.1894$ from the \Cloudy{} c25.00
emissivity file, a $1.4\%$ difference.

All 28 coefficient files carry updated provenance headers stating the CHIANTI
version, the population model, the energy convention, the fit range and the
maximum fit error, together with the density caveat of
Section~\ref{sec:limit}.

%==========================================================================
\section{Benchmarks after the fix}
\label{sec:bench}

Table~\ref{tab:bench} is the before-and-after on the two benchmark models.
HII40 moves from $9159$~K, above every published code, to $7894$~K, inside the
published range.  HII20 moves to $6376$~K, $0.4\%$ below the published floor.
These are the zero-density-limit (\code{low\_density}) temperatures, which the
density-suppressed default of Section~\ref{sec:limit} superseded one day later;
they are kept in the table because they are what the repair of the two defects
by itself is worth, and the audit trail of this memo is the point.  The
\emph{current} figures, in the last column of the table, are $8169$~K at HII40
and $6925$~K at HII20, and they change the HII20 verdict:

\begin{center}
\begin{tabular}{lcccl}
\toprule
 & published (nine codes) & \Cloudy{} c25.00 & \MoCHII{} & verdict \\
\midrule
HII40 & 7720--8199 & 8210 & \textbf{8169} & inside published; $0.50\%$ below \Cloudy{} \\
HII20 & 6402--6749 & 6910 & \textbf{6925} & above published; $0.22\%$ above \Cloudy{} \\
\bottomrule
\end{tabular}
\end{center}

\noindent
At HII20 it is not \MoCHII{} alone that exceeds the published range: \Cloudy{}
c25.00 exceeds it too, by $2.4\%$, and \MoCHII{} agrees with that
current-generation reference to $0.10\%$.  The published HII20 range is a
1995-era spread of nine codes, and both codes written since sit above it.  The
earlier reading of this section --- ``$0.4\%$ below the published floor'' ---
belonged to the zero-density fits and inverted the sign of the present
comparison; it is retained above only as the historical value.

\begin{table}[htbp]
\centering
\caption{The Lexington/Meudon benchmarks before and after the cooling fix, and
where they stand now.  ``published'' is the range over the nine codes of the
published compilation, which does not include C25; C25 is \Cloudy{} c25.00 run
on its own unmodified test-suite deck.  The $\nel \to 0$ column is the state
directly after the fit repair --- the zero-density (\code{low\_density})
cooling, threshold-aligned 32 ionizing bins, the explicit Monte Carlo diffuse
field with the \HeI{} channels, case-B recombination-line tables.  The
``current'' column is the production default: the density-suppressed
\code{local\_ne} cooling of Section~\ref{sec:limit}, continuous stellar photon
energies, Milne ground continua and the resolved \HeI{} diffuse cascade.  It
first raised $\langle T\rangle$ to 8211 and 7025~K and then settled at 8169 and
6925~K; Section~\ref{sec:limit} lists the intermediate steps.  The
recombination fraction was a diagnostic of the photon-budget argument and was
not re-measured for the current default.}
\label{tab:bench}
\begin{tabular}{lccccc}
\toprule
 & before & $\nel \to 0$ & current & C25 & published \\
\midrule
\multicolumn{6}{l}{\textbf{HII40}} \\
$L({\rm H}\beta)/10^{37}$ erg\,s$^{-1}$ & --- & 1.95 & 1.94 & 2.06 & 2.01--2.10 \\
$\langle T[N_{\rm p}N_{\rm e}]\rangle$ / K & 9159 & 7894 & \textbf{8169} & 8210 & 7720--8199 \\
$T_{\rm inner}$ / K                     & ---  & 7270 & 7535 & 7812 & 7370--8318 \\
$R_{\rm out}/10^{19}$ cm                & not predicted & 1.43 & 1.44 & 1.47 & 1.45--1.47 \\
$\langle {\rm He^+}\rangle/\langle {\rm H^+}\rangle$ & 0.667 & 0.715 & 0.706 & 0.741 & 0.715--0.829 \\
recombinations / $Q_{\rm H}$            & 0.888 & 0.953 & --- & --- & --- \\
\midrule
\multicolumn{6}{l}{\textbf{HII20}} \\
$L({\rm H}\beta)/10^{36}$ erg\,s$^{-1}$ & --- & 4.70 & 4.70 & 4.99 & 4.83--5.04 \\
$\langle T[N_{\rm p}N_{\rm e}]\rangle$ / K & 7370 & 6376 & \textbf{6925} & 6910 & 6402--6749 \\
$T_{\rm inner}$ / K                     & --- & 6752 & 6936 & 7155 & 6562--7224 \\
$R_{\rm out}/10^{18}$ cm                & --- & 8.69 & 8.89 & 8.97 & 8.70--9.01 \\
$\langle {\rm He^+}\rangle/\langle {\rm H^+}\rangle$ & 0.047 & 0.047 & 0.046 & 0.043 & 0.034--0.077 \\
recombinations / $Q_{\rm H}$            & --- & 0.971 & --- & --- & --- \\
\bottomrule
\end{tabular}
\end{table}

The photon budget closes at the same time, which is the confirming test of the
leakage argument of Section~\ref{sec:context}.  With the cooler gas the
ionization front no longer runs past the prescribed outer radius: the
recombination rate rises from $0.888$ to $0.953$ of $Q_{\rm H}$ at HII40, the
mean $x_{\rm H^+}$ in the outermost $2\%$ of the ionized volume falls to
$0.07$, and $R_{\rm out}$ becomes predictable at both benchmarks
($1.43\times10^{19}$ and $8.69\times10^{18}$~cm against published $1.45$--$1.47
\times10^{19}$ and $8.70$--$9.01\times10^{18}$).  $L({\rm H}\beta)$ follows to
within $3\%$ of the published values.

The line ratios move with the temperature.  Table~\ref{tab:lines} lists the
principal coolants; the systematic excess is gone, and the residual pattern is
what one expects from a model sitting at the cool edge of the range: the
temperature-sensitive optical lines of the doubly ionized species come out low
($[\SIII{}]$ 9532+9069 is the clearest case), while the infrared
fine-structure lines, which are nearly temperature-independent, sit inside the
published spread.

\begin{table}[htbp]
\centering
\caption{HII40 line intensities relative to ${\rm H}\beta$ after the fix.
``published'' is the range over the nine codes of the published compilation,
which does not include C25.
$[\OIII{}]$ 5007+4959 was $4.28$ before the fix.  Like the $\nel \to 0$ column
of Table~\ref{tab:bench}, these are the zero-density-limit values; the
density-suppressed default warms the model and raises the
temperature-sensitive optical lines.}
\label{tab:lines}
\begin{tabular}{lccc}
\toprule
Line & \MoCHII{} & C25 & published \\
\midrule
$[\OIII{}]$ 5007+4959       & 2.06   & 2.46   & 1.64--3.27 \\
$[\OIII{}]$ 52+88\,$\mu$m   & 2.41   & 2.40   & 2.10--2.42 \\
$[\OIII{}]$ 4363            & 0.0033 & 0.0046 & 0.0022--0.0070 \\
$[\OII{}]$ 3726+3729        & 1.90   & 2.32   & 1.60--2.26 \\
$[\NII{}]$ 6584+6548        & 0.626  & 0.669  & 0.669--0.852 \\
$[\NIII{}]$ 57.3\,$\mu$m    & 0.280  & 0.286  & 0.223--0.311 \\
$[\SIII{}]$ 9532+9069       & 0.91   & 1.00   & 1.13--1.44 \\
$[\SIII{}]$ 18.7\,$\mu$m    & 0.494  & 0.504  & 0.535--0.750 \\
$[\SII{}]$ 6716+6731        & 0.299  & 0.304  & 0.130--0.315 \\
$[\NeII{}]$ 12.8\,$\mu$m    & 0.202  & 0.200  & 0.177--0.217 \\
$[\NeIII{}]$ 15.5\,$\mu$m   & 0.278  & 0.216  & 0.263--0.429 \\
$\CIII{}]$ 1907+1909        & 0.029  & 0.062  & 0.041--0.091 \\
$[\SIV{}]$ 10.51\,$\mu$m    & 0.187  & 0.189  & --- \\
\bottomrule
\end{tabular}
\end{table}

%==========================================================================
\section{Density suppression of the line cooling at the local $\nel$}
\label{sec:limit}

The repaired fits are the $\nel \to 0$ limit, so they contain no collisional
de-excitation.  Lines whose critical density lies at or below the nebular
density are therefore overestimated.  On the same fixed \Cloudy{} state the
rebuilt Tier-1 line total is $1.116\times$ \Cloudy{}'s while the Tier-2 solve
at the local $\nel$ gives $1.011\times$; the difference is the density
suppression, and it is concentrated in exactly the expected places:

\begin{center}
\begin{tabular}{lccc}
\toprule
Channel & Tier-1 / \Cloudy{} & Tier-2 / \Cloudy{} & excess, points of $L_{\rm line}$ \\
\midrule
$[\CII{}]$ 158\,$\mu$m etc. & 2.763 & 0.940 & $+3.93$ \\
$[\OIII{}]$ (52 + 88\,$\mu$m share) & 1.126 & 1.024 & $+5.12$ \\
$[\NeIII{}]$ & 1.300 & 1.299 & $+0.73$ \\
$[\NIII{}]$  & 1.236 & 1.019 & $+0.56$ \\
$[\NII{}]$   & 1.095 & 0.990 & $+0.57$ \\
\midrule
line total & 1.116 & 1.011 & \\
\bottomrule
\end{tabular}
\end{center}

$[\CII{}]$ 158~$\mu$m has $\ncrit \sim 50$~cm$^{-3}$, so at $\nel = 110$ the
low-density limit overestimates it by $4.5\times$ (Table~\ref{tab:isolate},
$1.366\times10^{-21}$ against $3.048\times10^{-22}$).  The $[\OIII{}]$
far-infrared pair is the same effect at a higher $\ncrit$.  $[\NeIII{}]$ is
\emph{not} a density effect --- Tier-1 and Tier-2 agree with each other and
both exceed \Cloudy{} by $30\%$ --- so that row is a collision-data difference
between CHIANTI v11 and \Cloudy{}'s Ne\,\textsc{iii} data, and is noted here
as a separate open item.

Run with the zero-density fits used directly, the consequence for the
benchmarks was a slight cooling bias at $\nH = 100$: both models sat at the cool
edge of the published range, and HII20 fell $0.4\%$ below its floor.  \MoCHII{}
removes it by suppressing the line cooling to the local density.  A $(T, \nel)$ table,
built once at setup (\code{nlevel\_cooling\_mod}) by solving the same
$n$-level atom on a grid ($\log_{10}T = 2$--$6$, $\log_{10}\nel = -6$--$10$),
stores for each ion the suppression ratio
$S_{\rm sup}(T,\nel) = \Lambda_{\rm SE}(T,\nel)/\Lambda_{\rm SE}(T,\nel\!\to\!0)$
and the low-density reference; the runtime coefficient is
$\Lambda = \Lambda_{\rm fit}\,[\,f_{\rm cov}S_{\rm sup} + (1-f_{\rm cov})\,]$
with $f_{\rm cov} = \min(1,\ \Lambda_{\rm SE}(T,\nel\!\to\!0)/\Lambda_{\rm fit})$
the fraction of the low-density cooling the (often truncated) $n$-level model
resolves, so a truncated model cannot bias the total and $\nel \to 0$ returns
the fit exactly.  Because the table is built from the very solve that produces
the diagnostic lines, the thermal balance and the emergent line list are now
one physics.  This is the default (\code{par\%cooling\_model = local\_ne}); the
pure zero-density fits are kept as the \code{low\_density} option, exact in
diffuse gas and needed in any case for the planetary-nebula densities where the
suppression is much larger.

On the fixed \Cloudy{} state the default delivers the Tier-2 column of the
table above, $1.011\times$ \Cloudy{}'s line total in place of $1.116\times$,
and the converged benchmarks rose to \textbf{8211}~K at HII40 --- matching
\Cloudy{} c25.00's 8210~K to $0.7$~K --- and 7025~K at HII20 ($+1.7\%$ against
\Cloudy{}'s 6910~K), in place of the zero-density 7894 and 6376~K.  This
confirms the fixed-state prediction quantitatively: removing the far-infrared
overestimate --- the difference between the $1.116\times$ and $1.011\times$
totals --- is what carries the converged temperature the remaining distance to
the current-generation reference.  \NeIII{} remains the one channel still $30\%$
above \Cloudy{}, a collision-data difference rather than a density effect.

\paragraph{From 8211/7025 to the current 8169/6925~K.}  The two temperatures
above are those measured when \code{local\_ne} became the default; they are not
the current ones, and the difference is not in the cooling.  Three later
changes to the stellar and diffuse radiation field moved them, each recorded
with its own gate:

\begin{center}
\footnotesize
\begin{tabular}{llcc}
\toprule
Change & Plan document in \code{docs/} & HII40 & HII20 \\
\midrule
zero-density fits (\code{low\_density})   & this memo & 7894 & 6376 \\
\code{local\_ne} default                  & \code{COOLING\_LOCAL\_NE\_PLAN.md} & 8210 & 7024 \\
\CI{}/\NI{} $n$-level, energy order       & (same) & 8211 & 7025 \\
continuous stellar photon energies        & \code{CONTINUOUS\_ENERGY\_PHOTON\_PLAN.md} & 8194 & 6916 \\
Milne ground continua                     & \code{MILNE\_DIFFUSE\_LYC\_IMPLEMENTATION\_PLAN.md} & 8202 & 6922 \\
resolved \HeI{} diffuse cascade           & \code{HEI\_CASCADE\_DIFFUSE\_PLAN.md}, & \textbf{8169} & \textbf{6925} \\
                                          & \code{HEI\_584\_CONVERSION\_PLAN.md} & & \\
\bottomrule
\end{tabular}
\end{center}

\noindent
Sampling the stellar spectrum at continuous photon energies, rather than at
band-averaged energies, removed the discretization bias in the photoheating and
in the threshold cross sections.  Sampling the three ground recombination
continua from the Milne free-bound spectrum
$E^2\sigma_{\rm pi}(E)\exp[-(E-E_{\rm th})/kT]$ replaced an exponential
approximation, and the shift is above the Monte Carlo noise: the seed-to-seed
spread over four scrambled-Sobol ensembles is $0.002\%$ median on the line
ratios against a $0.40\%$ median Milne shift.  Resolving the \HeI{} diffuse
cascade replaced three constants --- statistical branch weights, a flat
$2^1S$ two-photon spectrum, and the 584~\AA{} photon treated as ionizing with
probability one --- by effective recombination coefficients and the
Drake--Victor--Dalgarno $2^1S$ spectrum, which lowers the effective ionizing
content of the helium cascade and so lowers HII40 by $40$~K.  None of this
touches the cooling function, and none of it disturbs the conclusion of this
memo: HII40 stays inside the published range throughout, and HII20 stays above
it --- within $1.7\%$ of \Cloudy{} c25.00 at the \code{local\_ne} default and
within $0.2\%$ from the continuous-energy step onward.

\paragraph{The H-collider term is still low-density (needs re-baselining).}
One line-cooling term is not yet density-suppressed: the H-impact $[\CII{}]$
158~$\mu$m / $[\OI{}]$ 63~$\mu$m fine-structure cooling of the
photodissociation-region package (\code{metal\_cooling\_H}), a two-level
low-density formula.  It is weighted by $x_{\rm HI}$, so it is $\sim$0 in the
ionized benchmark zones (the \HII{} gates are unaffected), but it needs the
same saturation in dense PDR gas ($\nH \sim 5\times10^3$ in the
\code{isrf\_cloud} example) --- a separate follow-up.  The photodissociation-region
test \code{tests/pdr/pdr\_L5.in} is otherwise unchanged: the PDR shell has
$\nel = A_{\rm C}\nH$ plus $6\%$ from S, Mg and Fe, $\Te = 258$~K,
$\CII{} = 1.000$, and there $\nel \sim 2\times10^{-2}$~cm$^{-3}$ is far below
$\ncrit$ so the low-density limit is the correct limit.  The stored PDR gate
references and the figures and examples built on the earlier cooling still need
refreshing to the density-suppressed default, and
\code{tests/g2a\_thermal/check\_g2a.py} still refers to a coefficient file
under its former name.

%==========================================================================
\section{Cooling inventories of the other Monte Carlo codes}
\label{sec:others}

The temperature of a photoionized nebula is set almost entirely by what is in
the cooling inventory, and Sections~\ref{sec:method}--\ref{sec:defects} showed
how a defect there hides itself behind the thermal balance.  The same audit was
therefore turned on the four Monte Carlo photoionization codes closest to
\MoCHII{} in method: the code of Wood et al.\ (2004), which is the WME column
of the benchmark tables; \CMI{}, its open successor, which inherits Wood's ion
ordering by its own account; \Mocassin{}; and \Torus{}, whose photoionization
loop borrows from both Wood's code and \Mocassin{}.  This is an internal record, so
everything found is written down, including defects on their side.  Line
numbers below were read from the sources in
\code{\textasciitilde/RT\_Codes/}.

\subsection{Wood et al.\ \code{ionize}}

\paragraph{Ion set.}  Ten five-level ions are read from \code{atom4.dat} in the
order N\,\textsc{i}, $\NII{}$, $\OI{}$, $\OII{}$, $\OIII{}$, $\NeIII{}$,
$\SII{}$, $\SIII{}$, $\CII{}$, $\CIII{}$ (\code{setcool.f}:25,
\code{DO 100 I=1,10}; \code{linecool1.f}:74), plus two two-level ions,
$\NIII{}$ and $\NeII{}$ (\code{linecool1.f}:160--171, abundances assigned at
\code{ioneng.f}:168--171).  The code is self-consistent, but two of its own
comments are not: the header list at \code{linecool1.f}:4--5 puts carbon in
slots 11--12 rather than 9--10, and the loop comment
\code{!omit 9,10 = C cooling} on line~74 is stale --- those slots are in the
loop.  \code{lines.f}:366--367 carries the corrected ordering.

\paragraph{Chlorine and argon data are present but unread.}
\code{atom4.dat} contains fourteen five-line blocks; blocks 11--14 are
Cl\,\textsc{ii}, Ar\,\textsc{iii}, Cl\,\textsc{iii} and Ar\,\textsc{iv}
(\code{atom4.dat}:51, 56, 61, 66).  They are never read, both because the read
loop stops at ten and because the arrays are dimensioned \code{(10,10)} and
\code{(10,5)} (\code{setcool.f}:7--8).  Chlorine and argon cooling are
therefore absent although their data sit in the file.

\paragraph{Defect: \NI{} and \OI{} line cooling.}
\code{linecool1.f}:42--64 evaluates hardcoded collision-strength fits into
\code{OM(1,*)} and \code{OM(3,*)} for \NI{} and \OI{}; then
\code{linecool1.f}:81--83, \emph{inside} the ion loop, overwrites the whole row,
\begin{center}
\code{DO MM=1,10 ; OM(J,MM)=CS(J,MM)*T4**CSE(J,MM) ; END DO}
\end{center}
so lines 42--64 are dead code.  The replacement values taken from
\code{atom4.dat} are zero: the \OI{} collision strengths and exponents
(\code{atom4.dat}:13--14) are all zero, and the \NI{} row
(\code{atom4.dat}:3) is zero except entries 5 and 10, which in the index order
of \code{setcool.f}:14 are the $(2,3)$ and $(4,5)$ transitions --- both between
excited levels.  Every ground-state excitation is zero, the level system
becomes homogeneous, and both ions cool exactly zero.  Reproducing
\code{linecool1.f}:85--123 with the file values at $T = 8000$~K,
$\nel = 100$~cm$^{-3}$ gives $\NI{}$ and $\OI{}$ cooling $= 0$ against
$4.75\times10^{-19}$~erg\,s$^{-1}$ for $\OIII{}$ on the same inputs.

There is a qualification that makes this worse rather than better.  \code{cs}
is one array shared by the cooling routine and the diagnostic routine, and
\code{lines.f}:406--428 writes its own \NI{}/\OI{} fits into \code{CS} rather
than \code{OM}; because the corresponding \code{CSE} entries are zero, those
values survive the overwrite at \code{linecool1.f}:81--83.  The diagnostic call
sits inside the iteration loop (\code{mcionize.f}:575--577) while the thermal
balance starts only at iteration 5 (\code{itbal.f}:27), so by the time
\code{linecool} runs, \code{CS(1,*)} and \code{CS(3,*)} hold values left from
the last cell of the previous iteration's diagnostic pass.  \NI{} and \OI{}
cooling is thus not identically zero at run time, but frozen at one unrelated
temperature, with no $T$ rescaling.  Either way it is not the physical term.

\paragraph{Defect: the cooling and diagnostic paths use different atomic data.}
The two two-level ions carry different collision strengths in the two routines:

\begin{center}
\begin{tabular}{lccc}
\toprule
 & cooling (\code{linecool1.f}) & diagnostic (\code{lines.f}) & CHIANTI v11 \\
\midrule
$\NIII{}$ 57.3\,$\mu$m $\Upsilon$ & 0.701\,$T_4^{0.136}$ (:31, :164) & 1.45 (:394, :540) & 1.378 \\
$\NeII{}$ 12.8\,$\mu$m $\Upsilon$ & 0.368 (:32, :169)                & 0.303 (:395, :545) & 0.314 \\
\bottomrule
\end{tabular}
\end{center}

\noindent
at $10^4$~K.  Only the cooling path carries the $T_4^{0.136}$ factor.  This one
can be adjudicated, because both codes describe the same transition: Wood's
$\NIII{}$ $A$-value, $4.77\times10^{-5}$~s$^{-1}$, is the CHIANTI v11 value to
three figures.  Evaluating the CHIANTI v11.0.2 collision strength through
\MoCHII{}'s own descaling gives $\Upsilon = 1.378$ at $10^4$~K and 1.296 at
8000~K, so the \emph{diagnostic} value 1.45 is right to $5\%$ and the
\emph{cooling} value is a factor $2.1$ too small.  $[\NIII{}]$ 57.3~$\mu$m is
$1.89\%$ of the HII40 cooling budget and lies below its critical density there,
so the cooling is low by nearly that factor in a channel of that size --- while
the code's own line output for the same transition uses the larger value.  For
$\NeII{}$, CHIANTI gives 0.314, again closer to the diagnostic value; the
cooling value is $17\%$ high.  The hardcoded \NI{} blocks are also assigned to
different transition indices in the two files (\code{linecool1.f}:44--50 versus
\code{lines.f}:408--414, with indices 3/4, 6/7 and 8/9 interchanged); which
ordering is intended cannot be settled from the sources, since neither cites
the row order of the table it transcribes.

\paragraph{\SIV{} slot never filled.}  The variable \code{csiv10} is declared
and passed through the argument lists (\code{lines.f}:46, 74, 204, 363, 388)
and used at \code{lines.f}:250 to fill \code{em(:,:,:,28)}, but no line of the
code ever assigns it.  There is no \SIV{} cooling anywhere in the code either.

\paragraph{\OIII{} $^1S_0$ level energy.}  \code{atom4.dat}:21 gives the
$\OIII{}$ levels as $0, 113.4, 306.8, 20271, 43072$~cm$^{-1}$.  The observed
$^1S_0$ energy is 43185.7~cm$^{-1}$ (CHIANTI \code{o\_3.elvlc}; \CMI{} carries
43186 at \code{LineCoolingData.cpp}:550), so the file value is low by
113.7~cm$^{-1}$, which is numerically the $^3P_1$ energy --- a transcription
slip.  It raises the ground-to-$^1S_0$ Boltzmann factor by $2.1\%$ at 8000~K
and shifts the nominal $[\OIII{}]$ 4363 wavelength to 4385.8~\AA{}, so the
line and any temperature diagnosed from it carry that bias.  The cooling itself
barely notices, because the $^1S_0$ channel is small.

\paragraph{H and He terms.}  The whole budget is
\code{loss = Lrec + Lff + Lc} (\code{ioneng.f}:191).  \code{RECcool}
(\code{glfunc.f}:15--29) has exactly two terms, H$^+$ and He$^+$; He$^{++}$ is
not tracked at all (\code{h0find.f} solves only $h^0$ and ${\rm he}^0$), so
there is no \HeIII{} recombination cooling and free-free is $Z = 1$ only.
There is no \HI{} collisional-excitation (Ly$\alpha$) cooling and no
collisional-ionization cooling.

\subsection{\CMI{}}

\paragraph{Ion set: thirteen ions, \SIV{} included.}  Ten five-level ions in
Wood's order (\code{LineCoolingData.hpp}:38--61 --- the file says outright that
``the order is historical: this was the order of the elements in Kenny's
code'') plus three two-level ions, $\NIII{}$, $\NeII{}$ and $\SIV{}$
(:68--77).  The \SIV{} data are real, not a placeholder: Martin, Zalubas \&
Musgrove (1990) energies, a Pradhan (1995) $A$-value of
$7.73\times10^{-3}$~s$^{-1}$, and fits to Saraph \& Storey (1999)
(\code{LineCoolingData.cpp}:1341--1376).  So \CMI{} carries the one channel
Wood omits entirely and \MoCHII{} was missing.

\paragraph{Atomic data provenance.}  The header
(\code{LineCoolingData.hpp}:126--133) states that many of the sources belong to
the IRON Project and were located through its online database, with the source
list at :136--182 (Berrington 1988; Froese Fischer \& Tachiev 2004; Galavis,
Mendoza \& Zeippen 1997, 1998; Kisielius et al.\ 2009; Lennon \& Burke 1994;
Tayal 2000, 2008; Tayal \& Zatsarinny 2010; Zatsarinny \& Tayal 2003).  The
compilation is a mixture, not exclusively IRON Project.

\paragraph{\NI{} and \OI{} are complete.}  Seven-parameter fits for all ten
transitions of each ion (\code{LineCoolingData.cpp}:112--118 for \NI{}, from
Tayal 2000; :353--359 for \OI{}, from Zatsarinny \& Tayal 2003 and Berrington
1988), fed abundances at \code{TemperatureCalculator.cpp}:374--390.  \CMI{}
does not carry Wood's \NI{}/\OI{} defect, nor the $\OIII{}$ $^1S_0$
transcription slip.

\paragraph{The He$^+$ difference is in the cooling, not the recombination
rate.}  The He$^+$ recombination \emph{coefficient} is identical in the two
codes: the same Verner \& Ferland (1996) HeIa fit with the same constants
$3.294\times10^{-11}$, $15.54$, $0.309$, $3.676\times10^{7}$, $1.691$
(\code{VernerRecombinationRates.cpp}:174--183 against Wood
\code{recomb.f}:17--19), so their ratio is unity at every temperature.  What
differs by about $11\%$ is the He$^+$ recombination \emph{cooling}
coefficient: Wood uses $L_{\rm He^+} = 2.6\,\nel n_{\rm He^+} T^{0.32}$
(\code{glfunc.f}:26) while \CMI{} uses $1.55\,\nel n_{\rm He^+} T^{0.3647}$
(\code{TemperatureCalculator.cpp}:489) --- which is exactly the expression Wood
carries commented out on the following line, \code{glfunc.f}:27.  Both are
attributed to Black (1981), Table~3.  The ratio Wood/\CMI{} is 1.146, 1.123,
1.111 and 1.077 at 5, 8, 10 and 20~kK.

\paragraph{Budget.}  The code's own comment at
\code{TemperatureCalculator.cpp}:347--351 lists three terms and no more: line
cooling, free-free, and H and He recombination.  There is no \HI{}
collisional-excitation cooling and no collisional-ionization cooling, and
He$^{++}$ is absent from the code altogether (\code{ElementNames.hpp} defines
only \code{ION\_H\_n} and \code{ION\_He\_n}), so free-free is $Z = 1$ only and
there is no \HeIII{} recombination cooling.

\subsection{\Mocassin{}}

\paragraph{Budget and solve.}  \code{coolInt = coolFF + coolRec + coolColl}
(\code{update\_mod.f90}:1251), with two dust terms added under
\code{lgDust .and. lgPhotoelectric} (:1253--1256).  \code{thermBalance} is
called from \code{iterateT} (:392), which re-solves the ionization balance at
every trial temperature (:365) and then steps on the residual
\code{heatInt - coolInt} (:397) --- the same arrangement as \MoCHII{}'s
bisection with the ionization solved inside it.  Collisionally excited line
photons are not transported: when a diffuse packet is drawn as a line photon it
escapes (\code{photon\_mod.f90}:925--927), so the line cooling is optically
thin, as in \MoCHII{}.  Only the resonance lines of a user list are followed,
and that is for dust absorption downstream of the cooling (\code{nResLines},
\code{common\_mod.f90}:533).

\paragraph{Ion set: whatever the data directory supplies.}  The heavy-element
sum runs \code{elem = 3, nElements} with \code{nElements = 30}
(\code{constants\_mod.f90}:54) and \code{ion = 1, min(elem+1, nstages)},
\code{nstages} defaulting to 7 and capped at 10
(\code{set\_input\_mod.f90}:82, :176--186).  An ion enters when its element has
a nonzero abundance and its file exists: \code{makeElements}
(\code{hydro\_mod.f90}:478) walks the 259 slots of \code{data/fileNames.dat}
(:511--518) and sets \code{lgDataAvailable} from the result (:701--731).  116
of the 259 files are present.  For the benchmark composition (He, C, N, O, Ne,
S) that is 26 ions at HII40's default \code{nstages} ---
C\,\textsc{i}--\textsc{iv}, N\,\textsc{i}--\textsc{v},
O\,\textsc{i}--\textsc{vi}, Ne\,\textsc{ii}--\textsc{vii} and
S\,\textsc{ii}--\textsc{vi} --- and 25 at the \code{nstages 6} of the HII20
deck (\code{benchmarks/gas/HII20/input.in}:22).  Adding an ion is a data
operation here too.  The coverage reaches higher ionization stages than
\MoCHII{}'s registry, which for these five elements stops at \CIII{}, \NIII{},
\OIII{}, \NeIII{} and \SIV{}; that is immaterial at 40~kK (\Cloudy{} puts
$0.007\%$ of the budget in $[\OIV{}]$) but not for the hard-spectrum models on
\MoCHII{}'s plan.

\paragraph{Level populations are a statistical-equilibrium solve at the local
$\nel$.}  \code{equilibrium} (\code{emission\_mod.f90}:1913--2122) interpolates
each $\Upsilon$ by a cubic spline in $\log_{10}T$ (:2015--2032), forms
$q_{lu} = 8.63\times10^{-6}\,\Upsilon\,e^{-\Delta E/kT}/(g_l\sqrt{T})$ from the
observed level energies of the file (:2040--2047), builds the rate matrix with
$\nel q$ in both directions plus the $A$-values (:2063--2089), solves it by LU
decomposition and normalizes the populations (:2091--2101), and returns
$A_{ji}\Delta E_{ji}n_j$ for every transition with a nonzero $A$
(:2104--2116).  Collisional de-excitation is therefore carried, and the
far-infrared fine-structure lines saturate above their critical density.
\MoCHII{} now does the same (Section~\ref{sec:limit}): where its $\nel \to 0$
fits deliver $1.116\times$ \Cloudy{}'s line cooling on the fixed \Cloudy{}
state, the $(T, \nel)$ table built from its own $n$-level solve delivers the
suppressed $1.011\times$ --- a difference that is a factor $4.5$ on $[\CII{}]$
158~$\mu$m alone --- and that suppressed value is what the thermal balance uses
by default.  Both codes therefore carry the density suppression in the thermal
balance.

The recombination-pumping term of the same routine is dormant.  It requires a
trailing \code{br, z, a\_r(4), a\_d(5)} block in the ion file
(\code{hydro\_mod.f90}:651--664, consumed at
\code{emission\_mod.f90}:1992--2004 and :2087); none of the 116 files carries
one, so \code{br = 0} and the populations are purely collisional --- the same
assumption as \MoCHII{}'s Tier-2 models.

Cooling and line output come from the same routine and the same files.  The
driver is duplicated, once inside the thermal balance
(\code{update\_mod.f90}:1264--1319, called at :1100) and once for the
emissivities (\code{emission\_mod.f90}:838--904), and the two differ only in
the unit factor ($1.9865\times10^{-16}$ against $1.9865\times10^{9}$).
\Mocassin{} therefore cannot carry the kind of cooling/diagnostic data split
found in Wood's code.

\paragraph{Two truncations.}  \code{nForLevels = 17}
(\code{constants\_mod.f90}:55) dimensions \code{forbiddenLines}
(\code{grid\_mod.f90}:89) and the line output uses all 17 levels, but the
cooling sum runs \code{j,k = 1..10} only (\code{update\_mod.f90}:1105--1106).
Thirty of the 116 ion models have more than ten levels, so for those the
cooling and the line list are not built from the same level set.  At nebular
temperatures this costs nothing: the $\OIII{}$ levels that fall outside the sum
start at 142382~cm$^{-1}$, a Boltzmann factor of $8\times10^{-12}$ at 8000~K.
Iron is the second case.  Its 142-level model (\code{nForLevelsLarge}, :56) is
written to a separate array, \code{forbiddenLinesLarge}
(\code{update\_mod.f90}:1278--1293), which the \code{coolColl} sum never reads,
so Fe\,\textsc{ii} line cooling is absent from the thermal balance although its
lines are output.  Both truncations are latent: every composition shipped with
the code has Fe = 0, and an iron- or calcium-bearing model would in any case
stop at \code{emission\_mod.f90}:1974--1978, since Fe\,\textsc{vi} (19 levels)
and Ca\,\textsc{vii} (27) exceed \code{nForLevels} and are reached at
\code{nstages} $\geq 6$ and $\geq 7$.

\paragraph{Atomic data: CHIANTI v7, with two pre-CHIANTI files among the main
coolants.}  Every ion file states its provenance in its own header.
\code{data/oiii.dat}:2 reads ``Atomic data from version 7 of the CHIANTI
database'', and so do all the benchmark files except two.
\code{data/siii.dat}:2--3 gives ``Upsilons from Mendoza (1983, Proc. IAU Symp.
103, p. 143); Aij's from Mendoza \& Zeippen (MNRAS, 199, 1025, 1982)'' on five
levels and four temperatures (5000--20000~K), and \code{data/ci.dat}:2--3 gives
Pequignot \& Aldrovandi (1976), with $A$-values from Nussbaumer \& Rusca
(1979), on six levels and six temperatures (500--20000~K).  $[\SIII{}]$ is
$15.2\%$ of the HII40 cooling budget (Table~\ref{tab:budget}), so \Mocassin{}'s
second-largest metal coolant runs on a 1983 compilation sampled at four
temperatures.

Table~\ref{tab:mocassinions} puts a number on what the vintage is worth.  For
each ion the $\nel \to 0$ line cooling was evaluated twice with the same
formula and the same upper levels, once from the \Mocassin{} file and once from
CHIANTI v11.0.2 through \MoCHII{}'s own descaling.  The procedure reproduces
the $\OIII{}$ cross-check quoted in Section~\ref{sec:context}: the total
$\Upsilon(^3P \to {}^1D)$ is 2.283 from the file against 2.262 from v11, a
ratio of 1.009.  The reading of the table is that the older data are
systematically \emph{higher} on three of the four largest metal coolants ---
$\SIII{}$ $1.223$ (its $\Upsilon(^3P \to {}^1D)$ is 8.39 against 6.96),
$\NII{}$ $1.199$, $\SII{}$ $1.163$, with $\NIII{}$ $1.154$ and $\OII{}$
$1.088$ --- while $\OIII{}$, $\NeIII{}$, $\SIV{}$ and $\OI{}$ agree to $3\%$.
$\CIII{}]$ is $1.457$ on a $0.41\%$ channel.

Two entries fall the other way, and the larger one is worth stating carefully.
\NI{} is a factor 17 low against v11, and this is a genuine change of
calculation rather than a transcription slip: \Mocassin{}'s file gives
$\Upsilon(^4S \to {}^2D_{5/2}) = 0.027$ at $10^4$~K and cites Tayal (2000),
\CMI{}'s independent fit to the same Tayal (2000) data gives $0.025$
(\code{LineCoolingData.cpp}:112--118 through the fit form at :1590--1598), and
CHIANTI v11, which uses Tayal (2006), gives $0.337$.  The two calculations
converge above $10^5$~K ($0.550$ against $0.563$).  So $\NI{}$ cooling and the
$[\NI{}]$ 5200 diagnostic differ by an order of magnitude between the 2000 and
the 2006 calculation, in exactly the partly neutral gas where \NI{} lives.
$\CI{}$ is $0.640$, the 1976 compilation against v11.

\begin{table}[htbp]
\centering
\footnotesize
\caption{\Mocassin{}'s ion models for the benchmark coolants: levels and
tabulated temperatures in \code{data/<ion>.dat}, the source of the effective
collision strengths as the file's own header gives it, and
$\Lambda_{\rm M}/\Lambda_{11}$, the ratio of the $\nel \to 0$ line cooling
computed from that file to the same quantity computed from CHIANTI v11.0.2
through \MoCHII{}'s descaling, at $T = 8000$~K with the same upper levels on
both sides.  All files are labeled CHIANTI v7 except \code{ci.dat} and
\code{siii.dat}, whose headers are quoted in the text.  The temperature grid is
the CHIANTI-wide $0.2$-dex grid from 100~K to $2.51\times10^{6}$~K where it has
23 points.}
\label{tab:mocassinions}
\begin{tabular}{llcclc}
\toprule
Ion & file & levels & $T$ pts & collision strengths & $\Lambda_{\rm M}/\Lambda_{11}$ \\
\midrule
$\CI{}$    & \code{ci.dat}    &  6 &  6 & Pequignot \& Aldrovandi (1976)      & 0.640 \\
$\CII{}$   & \code{cii.dat}   & 10 & 23 & Blum \& Pradhan (1992)              & 0.966 \\
$\CIII{}$  & \code{ciii.dat}  & 10 & 23 & Berrington et al.\ (1985, 1989)     & 1.457 \\
$\NI{}$    & \code{ni.dat}    & 13 & 23 & Tayal (2000)                        & 0.058 \\
$\NII{}$   & \code{nii.dat}   & 15 & 23 & Stafford et al.\ (1994)             & 1.199 \\
$\NIII{}$  & \code{niii.dat}  & 10 & 23 & Stafford, Bell \& Hibbert (1994)    & 1.154 \\
$\OI{}$    & \code{oi.dat}    &  7 & 23 & Zatsarinny \& Tayal (2003)          & 1.004 \\
$\OII{}$   & \code{oii.dat}   & 13 & 23 & Pradhan et al.\ (2006)              & 1.088 \\
$\OIII{}$  & \code{oiii.dat}  & 15 & 23 & Burke \& Lennon (1994)              & 1.028 \\
$\NeII{}$  & \code{neii.dat}  &  2 & 23 & Saraph \& Tully (1994)              & 0.905 \\
$\NeIII{}$ & \code{neiii.dat} &  9 & 23 & McLaughlin \& Bell (2000)           & 0.997 \\
$\SII{}$   & \code{sii.dat}   &  5 & 23 & Ramsbottom, Bell \& Stafford (1996) & 1.163 \\
$\SIII{}$  & \code{siii.dat}  &  5 &  4 & Mendoza (1983)                      & 1.223 \\
$\SIV{}$   & \code{siv.dat}   & 12 & 23 & Tayal (2000)                        & 0.992 \\
\bottomrule
\end{tabular}
\end{table}

\paragraph{\HI{}.}  Ly$\alpha$ and Ly$\beta$ collisional excitation from the
Mathis coefficients (\code{update\_mod.f90}:1072--1089).  The energy radiated
per excitation is taken as one Rydberg for both lines (\code{hcRyd}, :48)
rather than the 10.20 and 12.09~eV of the transitions, so the coefficients are
effective rather than literal --- and the product is right.  Against
\MoCHII{}'s CHIANTI v11 \HI{} low-density fit
(\code{data/atomic/cooling\_h\_1.txt}) the expression gives
$2.211\times10^{-25}$ against $2.276\times10^{-25}$ at 8000~K and
$4.140\times10^{-24}$ against $4.223\times10^{-24}$ at $10^4$~K, that is
$0.97$--$0.98$, falling to $0.94$--$0.95$ at 6 and 20~kK; with the literal line
energies it would be $25\%$ low, because the two-term form has to stand in for
the whole Lyman series.  The hard cut \code{coolColl = 0} below 5000~K (:1079)
removes nothing measurable.  There is no separate two-photon or free-bound
cooling term; free-bound sits inside \code{coolRec}.

\paragraph{He.}  There is no \HeI{} or \HeII{} collisional-excitation cooling:
the collisionally excited line sum starts at \code{elem = 3} in both drivers
(\code{update\_mod.f90}:1103, :1272; \code{emission\_mod.f90}:846).  He$^{++}$
free-free and recombination cooling are present (:1014--1035), the Hummer
(1994) Table~1 fits with the polynomial doubled and evaluated at
$\log_{10}(T/4)$, and He$^{+}$ free-free has its own fit (:1047--1048,
attributed to Hummer \& Storey 1998, Table~6).  The He$^{+}$ \emph{recombination}
cooling, however, is the H$^{+}$ polynomial copied coefficient for coefficient:
compare :1059--1062 with :989--991, all five constants identical, while the
comment above it cites the helium source.  At $10^4$~K it gives
$2.40\times10^{-25}$~erg\,cm$^3$\,s$^{-1}$ against $3.88\times10^{-25}$ for
\MoCHII{}'s $k_{\rm B}T\alpha_{\rm B}({\rm He\,\textsc{ii}})$,
$4.95\times10^{-25}$ for Wood's Black-1981 coefficient and
$4.46\times10^{-25}$ for \CMI{}'s, so it is low by a factor $1.6$--$2.1$
against all three.  Only part of that is the copy:
$\alpha_{\rm B}({\rm He\,\textsc{ii}})$ exceeds
$\alpha_{\rm B}({\rm H\,\textsc{ii}})$ by $8\%$ at this temperature, so the
substitution itself is a small error and the remainder is which mean energy per
recombination the various coefficients represent.  The channel is worth of
order $1\%$ of the HII40 budget.

\paragraph{Recombination and free-free elsewhere agree well.}  \Mocassin{}'s
H$^{+}$ recombination cooling equals \MoCHII{}'s Hui \& Gnedin \emph{case-B}
coefficient to $0.1\%$ at 8000~K and $1.1\%$ at $10^4$~K ($2.346$ against
$2.343\times10^{-25}$, and $2.403$ against $2.376\times10^{-25}$), and its
He$^{++}$ term is $0.97\times$ the case-B \HeIII{} coefficient, so its
recombination cooling is case-B bookkeeping.  Its free-free fits correspond to
thermally averaged Gaunt factors of $1.259$ and $1.272$ at $Z = 1$ (8000 and
$10^4$~K) and $1.179$, $1.191$ at $Z = 2$, against $1.252$, $1.266$ and
$1.176$, $1.187$ from the van Hoof et al.\ (2014) table \MoCHII{} can select:
agreement better than $1\%$.  What is missing is the metals.  Free-free is
summed over H$^{+}$, He$^{+}$ and He$^{++}$ only, with no net-charge term for
the metal ions (\MoCHII{} adds one, worth $\sim2\times10^{-4}$ of the H/He
term), and there is no metal recombination cooling ($0.05\%$ of the \Cloudy{}
budget, Table~\ref{tab:budget}; \MoCHII{} omits it too).

\paragraph{Collisional ionization is not in the budget.}  \code{collIon}
(:1378, Drake \& Ulrich 1980) is computed for hydrogen only, explicitly zeroed
for every other species (:1485), and enters the ionization ratios (:1490,
:1496).  It never reaches \code{coolInt}.

\paragraph{Dust.}  \code{setDustGasCollHeatCool} (:811--931) is the gas--grain
collisional exchange: ions (:831--892) and electrons (:894--924) striking
grains at the local grain potential, accumulated as \code{gasDustColl\_g} for
the gas and \code{gasDustColl\_d} for the grains.  It reaches the budget
together with the photoelectric heating at :1253--1256, and both are computed
only under \code{lgDust .and. lgGas .and. lgPhotoelectric} (:368--386), so a
dusty model run with the photoelectric switch off has neither.  This is the one
cooling channel \Mocassin{} has and \MoCHII{} does not.  Its size at benchmark
conditions is negligible: with the standard gas--grain rate (Hollenbach \&
McKee 1989) at $\nH = 100$~cm$^{-3}$, $T = 8000$~K, $T_{\rm d} = 50$~K and
Milky Way dust it is $\sim5\times10^{-24}$~erg\,cm$^{-3}$\,s$^{-1}$ against
$\sim4\times10^{-20}$ for the lines, of order $10^{-4}$.  It scales as $\nH^2$
like the lines, so it matters only at planetary-nebula and
photodissociation-region densities, and both benchmark models are dust-free.
\MoCHII{} routes grain energy to the infrared instead (\code{Heat\_dust}) and
carries grain photoelectric heating with grain recombination cooling
(\code{par\%grain\_pe}).

\paragraph{Net comparison with \MoCHII{}.}  The dominant difference is now
the collision data.  Both codes suppress the metal line cooling to the local
$\nel$ in the thermal balance --- \Mocassin{} by solving the level populations
there, \MoCHII{} through the $(T, \nel)$ table built from its own $n$-level
solve (Section~\ref{sec:limit}) --- so the $\sim$$10\%$ that the suppression
removes at $\nH = 100$~cm$^{-3}$ is common to both.  What remains is that
\Mocassin{}'s collision strengths are $15$--$22\%$ higher than the current ones
on three of the four largest metal coolants, which adds cooling, while
\MoCHII{} carries the current CHIANTI~v11 strengths.  Neither code's temperature
can therefore be attributed to a single data difference.  On channels rather
than data, \MoCHII{} carries three that \Mocassin{} does not ---
collisional-ionization cooling, metal free-free with the net ion charge, and
H-impact fine-structure cooling for the photodissociation-region path
(\Mocassin{}'s $\Upsilon$ tables are electron-impact only) --- and \Mocassin{}
carries one that \MoCHII{} does not, the gas--grain collisional term.

\subsection{\Torus{}}

\Torus{} (Harries et al.\ 2019) is a Monte Carlo radiation transport and
hydrodynamics code, and its photoionization loop --- the reference the code
registers for itself when that loop runs is Haworth \& Harries (2012),
\code{citations\_mod.f90}:52 --- is a full photoionization-equilibrium plus
thermal-balance calculation on the same kind of octree \MoCHII{} uses, so it
belongs in this comparison.  Its provenance runs through both of the codes
above.  The free-free and photoheating expressions carry the comments ``Kenny
equation 22'', ``Kenny equation 23'' and ``equation 21 of kenny's''
(\code{photoionAMR\_mod.F90}:6465--6467, :7813), that is Wood's paper; the
\HI{} collisional-excitation block is \Mocassin{}'s character for character
(below); and two of the ion data blocks say in their own comments that they
were tuned to agree with \Mocassin{} and with Ercolano's $\SIII{}$ data
(\code{ion\_mod.F90}:609, :1138).

\paragraph{Two copies of the module; the MPI one is the production path.}
\code{photoion\_mod.F90} (3406 lines) is the serial loop, reached from
\code{physics\_mod.F90}:412; \code{photoionAMR\_mod.F90} (11179 lines) is the
MPI/AMR loop, reached from \code{physics\_mod.F90}:437 and compiled only under
\code{\#ifdef MPI}.  The two carry independent copies of \code{HHeCooling},
\code{getHeating}, \code{metalcoolingRate} and \code{solveIonizationBalance},
and they are not identical: the serial one always adds the gas--grain term and
has no dust radiative term (\code{photoion\_mod.F90}:1368), the MPI one makes
both conditional.  Everything quoted below is the MPI version unless stated.

\paragraph{Method: Lucy estimator with immediate re-emission.}  The packet
energy is $\epsilon/\Delta t = L_{\rm ion}/N_{\rm packet}$ (:3626), and the
photoionization rate of every ion is
the Lucy (1999) path-length estimator, $(\epsilon/\Delta t/V)\sum \ell\,
\sigma_i/h\nu$ (accumulated in \code{updateGrid}:6691--6692, consumed at
:6880); the photoheating uses the same estimator with the weight
$\ell\,\sigma\,(h\nu - h\nu_{\rm th})/h\nu$ (:6717--6728).  When a packet is
absorbed it is \emph{immediately re-emitted}: the cell's diffuse emissivity
spectrum is assembled once per radiation step from the Lyman continua
(:3325), the higher continua (:3326), the \HI{} recombination lines (:3328),
the collisionally excited metal lines (:3330) and the dust continuum (:3335),
a new frequency is drawn from it (:3372) and the packet continues in a random
direction with its energy unchanged (:3374).  So \Torus{} does not let its
collisionally excited line photons escape the way \Mocassin{} and \MoCHII{}
do: they re-enter the transport, and in a photoionized nebula their only sink
is dust absorption, so this is how line-cooling energy reaches the infrared.
Because the re-emitted packet keeps the absorbed energy and takes only its
\emph{shape} from the emissivity, the emergent line luminosities equal the
cooling-rate line luminosities only in a converged model, where heating equals
cooling by construction of the temperature solve.  An on-the-spot option
exists: \code{nodiffuse} zeroes everything above 13.6~eV in the spectrum
before the redraw (:3347--3349, \code{removeLymanSpectrum} :7858--7868).

\paragraph{The temperature: three modes.}  In the equilibrium mode the
ionization balance and the temperature alternate in an inner loop per octal
until both change by less than $1\%$ (:3908--3927), and
\code{calculateThermalBalanceNew} (:6109--6286) finds the root of
$\Gamma_{\rm tot} - \Lambda_{\rm tot}$ by Brent's method, bracketing first on
$0.8T$ and $1.2T$ and falling back to $[T_{\rm min}, T_{\rm max}]$
(:6165--6176), then under-relaxing the step by $0.6$, or by $0.1$ after
iteration~10 (:6132--6135).  Unlike \Mocassin{} and \MoCHII{}, the ionization
balance is \emph{not} re-solved at each trial temperature --- the call that
would do it sits commented out inside the cooling function
(\code{photoion\_mod.F90}:1354) --- so the coupling is carried by the outer
alternation instead.  The time-dependent mode integrates
$de/dt = \Gamma - \Lambda$ in ten substeps with the same cooling function
(\code{calculateThermalBalanceTimeDep} :6364--6411).  A third mode,
\code{quickthermal} (\code{inputs\_modV2.F90}:3948), skips the balance
altogether and sets $T = T_{\rm min} + (10^4~{\rm K} - T_{\rm min})\,
x_{\rm H^+}$ (:5974--5990).  Cells the Monte Carlo has undersampled are forced
to $T_{\rm min}$ (:6279--6280), which for the default $T_{\rm min} = 3$~K
(\code{inputs\_modV2.F90}:65) is a placeholder rather than a physical
temperature.

\paragraph{Budget: seven terms, and the code counts them itself.}
\code{HHeCooling} (:6413--6648) accumulates, in order: free-free (:6465--6467),
\HI{} collisional excitation (:6500--6520), H$^+$ recombination cooling
(:6571), He$^+$ recombination cooling (:6582), metal line cooling (:6593),
the gas--grain collisional term (:6604--6612) and dust radiative cooling
(:6620--6628).  The routine's own diagnostic branch asserts
\code{nRates /= 7} (:6636), so seven is the intended count.

Two expressions in the same routine are dead.  A Hummer (1994) free-free
coefficient and a Hummer (1994) H$^+$ recombination-cooling coefficient are
evaluated at :6480--6496 and accumulated into \code{becool}, which is used
only by a debug \code{write} (:6528--6534) and never reaches
\code{coolingRate}.  The free-free that is actually used is the ``Kenny''
form, $1.42\times10^{-27}(n_{\rm H^+} + n_{\rm He^+})\,\nel\,
g_{\rm ff}\sqrt{T}$ with $g_{\rm ff} = 1.1 + 0.34\exp[-(5.5-\log_{10}T)^2/3]$.
That $g_{\rm ff}$ is $1.245$ at 8000~K and $1.261$ at $10^4$~K, against
$1.252$ and $1.266$ from the van Hoof et al.\ (2014) table \MoCHII{} can
select --- agreement to $0.6\%$.  The sum runs over H$^+$ and He$^+$ only, both
at $Z = 1$: He$^{++}$ is tracked in the ionization balance
(\code{ion\_mod.F90}:215) but contributes neither $Z = 2$ free-free nor
recombination cooling, and the metals contribute no free-free.

\paragraph{\HI{} collisional excitation is \Mocassin{}'s expression.}
:6502--6518 carries \code{ch12 = 2.47e-8}, \code{ch13 = 1.32e-8},
\code{ex12 = -0.228}, \code{ex13 = -0.460}, \code{th12 = 118338.},
\code{th13 = 140252.}, the same \code{hcRyd} per excitation and the same hard
cut below 5000~K as \code{update\_mod.f90}:1071--1089, under the same comment
``Mathis, Ly alpha, beta''.  It therefore reproduces the \Mocassin{} numbers
quoted above exactly: $2.211\times10^{-25}$ at 8000~K and
$4.140\times10^{-24}$ at $10^4$~K, that is $0.97$--$0.98$ of \MoCHII{}'s
CHIANTI~v11 \HI{} low-density fit.

\paragraph{Recombination cooling: two tabulated coefficients, no cited source.}
$\nel n_{\rm H^+}kT\beta_{\rm H}$ and $\nel n_{\rm He^+}kT\beta_{\rm He}$ with
$\sqrt{T}\beta$ read from a 31-point table over $\log_{10}T = 1$--7 for
hydrogen and an 18-point table over 1--4.4 for helium, linearly interpolated in
$\log_{10}T$ and set to zero above the helium table's end (:6426--6435,
:6544--6569).  Neither table carries a provenance comment.  Numerically the
hydrogen one is case-B recombination cooling: $kT\beta_{\rm H} =
2.360\times10^{-25}$ at 8000~K and $2.397\times10^{-25}$ at $10^4$~K, which is
$1.007$ and $1.009$ times \MoCHII{}'s Hui \& Gnedin case-B coefficient and
$1.006$ and $0.997$ times \Mocassin{}'s.  The helium table is a genuinely
separate table, not the hydrogen one reused --- which is the defect recorded
on \Mocassin{}'s side above --- and it gives $2.619\times10^{-25}$ at 8000~K
and $2.706\times10^{-25}$ at $10^4$~K.  At $10^4$~K that is $1.13$ times
\Mocassin{}'s copied value and $0.70$ times \MoCHII{}'s
$kT\alpha_{\rm B}({\rm He\,\textsc{ii}})$, so it sits at the low end of the
spread but not at \Mocassin{}'s extreme.

\paragraph{Dust enters as two terms.}  With
\code{decouplegasdust} false, which is the default
(\code{inputs\_modV2.F90}:3951), the gas and the grains share one temperature,
the grain thermal emission $4\pi\kappa_{\rm P}(\sigma/\pi)T^4$ is added to the
cooling side of the same balance (:6620--6628) and the grain heating rate is
added to the heating side (:7818--7823).  With it true the two temperatures are
separate, the gas--grain collisional exchange
$n_{\rm gas}n_{\rm gr}\sigma_{\rm gr}v_{\rm p}f\,2k(T_{\rm gas}-T_{\rm d})$
becomes a gas cooling term (:6604--6612;
\code{photoion\_utils\_mod.F90}:1634--1680, with accommodation factor
$f = 11.8$ in ionized gas and $0.16$ or $1$ in neutral gas, :1666--1674), and
the grain temperature is solved from radiative equilibrium plus that exchange
(:6253--6271).  So \Torus{} carries the gas--grain collisional channel that
\Mocassin{} has and \MoCHII{} does not, as an option, and it additionally
routes grain emission through the gas thermal balance, which no other code
here does.

\paragraph{Level populations are a statistical-equilibrium solve at the local
$\nel$.}  \code{solvePops} (\code{photoion\_utils\_mod.F90}:923--997) builds
the rate matrix with $\nel q$ in both directions plus the $A$-values
(:954--974), replaces the first row by the normalization condition
(:942--943), solves it by LU decomposition (:986) and normalizes; the
collisional rates are $q = 8.63\times10^{-6}\,\Upsilon/(g\sqrt{T})$ with the
Boltzmann factor from the observed level energies (\code{getCollisionalRates}
:999--1046), $\Upsilon$ linearly interpolated in $T$ and clamped at the table
ends with a one-time warning (:1015--1027).  So \Torus{}, like \Mocassin{} and
like \MoCHII{}'s density-suppressed $\Lambda(T,\nel)$ (Section~\ref{sec:limit}),
carries collisional de-excitation inside the thermal balance and saturates the
far-infrared fine-structure lines above their critical density.  Two
qualifications.  Every caller passes a population array dimensioned
\code{pops(10)} (\code{photoionAMR\_mod.F90}:7712, :7732, :8089;
\code{photoion\_utils\_mod.F90}:1115), so no ion model may exceed ten levels;
the largest in the file is $\CII{}$ with exactly ten.  And the
recombination-pumping term is dormant here too: \code{getRecombs} is called
(:940) but the line that would put its rates into the right-hand side is
commented out and \code{matrixB(2:n)} is set to zero (:945), so the populations
are purely collisional --- the same assumption as \MoCHII{}'s Tier-2 models.

\paragraph{Ion set: fifteen line-cooling ions, hardcoded.}  There are no data
files.  \code{addIons} (\code{ion\_mod.F90}:194--388) builds the registry ---
\HI{}, \HII{}, \HeI{}, \HeII{}, \HeIII{}, then C\,\textsc{i}--\textsc{iv},
N\,\textsc{i}--\textsc{iii}, O\,\textsc{i}--\textsc{iii},
Ne\,\textsc{i}--\textsc{iii} and S\,\textsc{i}--\textsc{iv} under
\code{usemetals}, twenty-two ions in all --- and
\code{addLevels} (:428--587) and \code{addTransitions} (:603--1186) fill levels
and transitions from \code{select case} blocks on the species name, the level
energies taken from NIST by the file's own statement (:433).  Fifteen ions end
up with transitions and therefore with cooling: C\,\textsc{i}--\textsc{iv},
N\,\textsc{i}--\textsc{iii}, O\,\textsc{i}--\textsc{iii}, $\NeII{}$,
$\NeIII{}$ and S\,\textsc{ii}--\textsc{iv}.  Ne\,\textsc{i} and
S\,\textsc{i} are in the ionization balance but have neither levels nor lines.
$\NI{}$ and $\OI{}$ are complete --- ten transitions each with tabulated
$\Upsilon(T)$, the full set for their five levels (:812--847, :914--937) ---
and $\SIV{}$ is present as a two-level ion with
$A = 7.74\times10^{-3}$~s$^{-1}$ (:1179), which is \CMI{}'s Pradhan (1995)
value to three figures.  So \Torus{} carries neither of the two channels
Wood's code loses, and like \CMI{} it has the $\SIV{}$ channel \MoCHII{} was
missing.  There is no chlorine or argon line cooling: argon ions exist only
under \code{xraymetals} and have no level data at all.

The transition energy is the difference of the two observed level energies
(\code{addTransition2} :1266), the stated wavelength is kept separately for
line identification only, and a warning fires if the two disagree by more than
$5\%$ (:1273--1279).  \Torus{} therefore cannot carry \MoCHII{}'s D2 defect by
construction.  Cooling and line output run through the same routine and the
same data --- \code{metalcoolingRate}:7742 and
\code{cellLineLuminosity}:7707--7722 are the same expression with the same
\code{solvePops} call --- so it cannot carry Wood's cooling/diagnostic split
either.

\paragraph{Collision-strength vintage: mixed, and \NI{} is the modern
calculation.}  The tables are per transition and per ion, with temperature
grids from two points ($\CIV{}$) to sixteen ($\SII{}$, $\SIII{}$), and their
lineage is visible in the comments: carbon was ``freshly mod[ded] \dots\ to
agree with MOCASSIN'' (:609) and the live $\SIII{}$ block is the fourth
attempt, labeled ``Matching Barbara Ercolano's S III data as closely as
possible'' (:1138), with a Tayal \& Gupta (1999) block and two others left
commented out above it (:1060--1136).  Two spot checks against the numbers
above.  For $\OIII{}$ the total $\Upsilon(^3P \to {}^1D)$ from :967--991 is
$2.119$ at 8000~K and $2.184$ at $10^4$~K, so $0.965$ of CHIANTI~v11's $2.262$
and $0.957$ of \Mocassin{}'s v7 file value $2.283$ --- the same few-percent
band all three sit in.  For $\NI{}$, however,
$\Upsilon(^4S \to {}^2D_{5/2})$ is $0.2338$ at 8000~K and $0.2877$ at $10^4$~K
(:817--819), which is $0.85$ of CHIANTI~v11's $0.337$ and an order of magnitude
above \Mocassin{}'s $0.027$ and \CMI{}'s $0.025$.  \Torus{}'s nitrogen is thus
on the modern side of the Tayal (2000) versus Tayal (2006) split, and its
$\NI{}$ cooling and $[\NI{}]$ 5200 diagnostic differ from theirs by that
factor.

\paragraph{Defect: a Fortran exponent typo makes an $\SII{}$ $A$-value
negative.}  In the \code{addTransition2} call for the $\SII{}$
$^2P^{\rm o}_{3/2} \to {}^2P^{\rm o}_{1/2}$ fine-structure transition
(\code{ion\_mod.F90}:1055) the $A$-value argument is written \code{1.03-6},
which in Fortran is not $1.03\times10^{-6}$ but the expression
$1.03 - 6 = -4.97$.  The transition therefore enters with
$A = -4.97$~s$^{-1}$, which exceeds the sum of the three genuine decay rates
out of that level ($0.537$~s$^{-1}$) and makes its total radiative loss rate
negative.  Reproducing \code{solvePops} with the file's own five levels, ten
transitions and $\Upsilon(T)$ tables gives a \emph{negative}
$^2P^{\rm o}_{3/2}$ population at every temperature and density tried
($-1.08\times10^{-8}$ at $T = 8000$~K, $\nel = 100$~cm$^{-3}$, falling to
$-2.15\times10^{-5}$ at 20~kK and $10^4$~cm$^{-3}$), which the
\code{pops = max(0, \dots)} clamp at \code{photoion\_utils\_mod.F90}:992 then
sets to zero.  The consequences are specific.  The three $[\SII{}]$ lines whose
upper level is $^2P^{\rm o}_{3/2}$ --- 4068.6, 10286.7 and 10320.4~\AA{} ---
emit exactly zero, while the three whose upper level is
$^2P^{\rm o}_{1/2}$ --- 4076.4, 10336.3 and 10373.3~\AA{} --- come out a factor
$3.2$ too strong, because the negative level acts as a sink that inflates
$^2P^{\rm o}_{1/2}$ from $7.28\times10^{-8}$ to $2.34\times10^{-7}$.  The
intrinsic 4069/4076 ratio, $3.04$ on the intended $A$, comes out $0$.  The
\emph{total} $\SII{}$ cooling barely notices, because the two errors have
opposite sign: $1.005$ of the intended value at 8000~K and $\nel = 100$,
$1.014$ at $\nel = 10^4$, $1.025$ at 20~kK and $10^4$.  With $[\SII{}]$ at
$2.16\%$ of the HII40 budget (Table~\ref{tab:budget}) the thermal balance is
unaffected at the $10^{-4}$ level; what is destroyed is the $[\SII{}]$ auroral
diagnostic.  The ``Negative contribution from'' warning writer
(\code{photoionAMR\_mod.F90}:7745--7749) exists but tests the summed ion rate,
which stays positive.

\paragraph{Defect: the $\NII{}$ $^1S_0$ statistical weight.}
\code{ion\_mod.F90}:479 enters the $\NII{}$ $^1S_0$ level with $J = 2$, so
$g = 5$ instead of $1$; every other $^1S_0$ in the file is entered with
$J = 0$ ($\CIII{}$ :456, $\OIII{}$ :516, $\NeIII{}$ :532, $\SIII{}$ :549).
Only the de-excitation rate divides by the upper level's $g$
(\code{photoion\_utils\_mod.F90}:1042), so the error is confined to
collisional de-excitation out of $^1S_0$; with
$\ncrit \sim 2\times10^{7}$~cm$^{-3}$ for that level it is negligible at
nebular densities and wrong by a
factor $5$ in the de-excitation rate wherever it is not.  Mg\,\textsc{i}
carries the same slip on its $^1P^{\rm o}_1$ level (:564), harmlessly, because
magnesium has no transitions.

\paragraph{The X-ray metal set cannot be used as written.}  \code{xraymetals}
adds twenty-one Mg, Si, Ar and Fe ions to the registry
(\code{ion\_mod.F90}:296--375), but \code{returnAbundance} (:1317--1345) has no
case for $Z = 12$, 14, 18 or 26 and returns \code{tiny(a)} with a runtime
message, and none of the four elements has transition data (magnesium has
levels but no transitions; silicon, argon and iron have neither).  So the
switch contributes no cooling and no opacity of consequence.  A latent
companion: Mg\,\textsc{ii} declares two levels both named \code{"2S\_1/2"}
(:567, :570) and \code{returnLevel} (:1298--1315) returns the first match, so
the second would be unreachable if transitions were ever added.

\paragraph{What is absent.}  There is no collisional ionization anywhere:
\code{solveIonizationBalance} (:6838--6918) forms each stage ratio as
(photoionization $+$ charge-exchange ionization) over (radiative $+$
dielectronic recombination $+$ charge-exchange recombination), with no
collisional term, so there is no collisional-ionization cooling either.  There
is no \HeI{} or \HeII{} collisional-excitation cooling --- the metal sum starts
at registry entry~5 (:7737), which is \HeIII{} with no level data, so
effectively at $\CI{}$, and helium has no level data at all.  There is no
\HeIII{} recombination cooling, no $Z = 2$ free-free, no metal free-free, no
metal recombination cooling and no two-photon term (the two-photon routine is
commented out, :7775--7799).  On the heating side only \HI{} and \HeI{}
photoionization are counted (:6717--6728, :7812--7817): \HeII{} photoionization
heating is missing although He$^{++}$ is tracked, and there is no metal
photoheating, no grain photoelectric heating and no cosmic-ray heating in the
photoionization loop.  A separate photodissociation-region module does carry
Weingartner \& Draine photoelectric heating, cosmic-ray, turbulent, chemical,
gas--grain and H$_2$-formation heating (\code{pdr\_utils\_mod.F90}:366--800),
invoked from the photoionization loop under \code{pdrcalc}
(\code{photoionAMR\_mod.F90}:527, :997), but that is a second thermal balance
with its own inventory and is outside this comparison.  One point in
\Torus{}'s favor on the electron side: metal electrons are always in $\nel$,
since \code{returnNe} (\code{ion\_mod.F90}:1410--1433) sums $Z - N$ over the
whole registry, where \MoCHII{} needs \code{par\%metal\_ne}.

\paragraph{The benchmark directory.}  A Lexington HII40 model is shipped with
the code.  \code{benchmarks/HII\_region/} and \code{benchmarks/HII\_regionMPI/}
are the same model in the OpenMP and MPI builds: \code{geometry lexington}, a 1D octree
at depth 8, one blackbody source at $T_{\rm eff} = 40000$~K with
$R = 18.67\,R_\odot$ (\code{parameters.dat}), and $\nH = 100$~cm$^{-3}$ uniform
outside \code{rinner} (\code{calcLexington}, \code{amr\_mod.F90}:7557--7612).
The default abundances --- He $0.1$, C $2.2\times10^{-4}$, N $4\times10^{-5}$,
O $3.3\times10^{-4}$, Ne $5\times10^{-5}$, S $9\times10^{-6}$
(\code{inputs\_modV2.F90}:4107--4133) --- are exactly the Lexington/Meudon
HII40 composition, so \Torus{}'s photoionization defaults are set by the same
benchmark \MoCHII{} is measured against.  The output dump
(\code{dumpLexingtonMPI}, \code{photoionAMR\_mod.F90}:7353--7480) writes 500
radial points from 1 to 8~pc: radius, $\Te$, then $\log x$ for \HI{}, \HeI{},
$\OII{}$, $\OIII{}$, $\CII{}$, $\CIII{}$, $\CIV{}$, $\NII{}$, $\NIII{}$,
N\,\textsc{iv}, Ne\,\textsc{i}, $\NeII{}$, $\NeIII{}$, Ne\,\textsc{iv} and
$T_{\rm d}$, plus \code{ne.dat} and an \code{orates.dat} carrying the cell
heating and cooling rates.

The reference file \code{lexington\_benchmark.dat} is committed to the
repository, so it is \Torus{}'s own stored regression output rather than a
published table, and \code{comparelex.f90} checks a new run against it in
$\Te$ and one ion at a $10\%$ tolerance.  It is still usable as a fifth
column: weighting the stored profile the way the published benchmark tables do
gives $\langle T[N_{\rm p}N_{\rm e}]\rangle = 8196$~K, with the
$x_{\rm HI} = 0.5$ front at $4.75$~pc.  Two caveats before that number is
used.  The dump starts at 1~pc, so the innermost parsec is not in the average;
and it is a single stored run at whatever packet and iteration count produced
it, not a converged reference in the sense of Section~\ref{sec:bench}.  Taken
at face value it sits at the top of the $7720$--$8199$~K published range of
Table~\ref{tab:bench}, $0.4\%$ above \MoCHII{}'s current $8169$~K and $0.2\%$
below \Cloudy{} c25.00's $8210$~K (and above \MoCHII{}'s zero-density
$7894$~K).  \code{benchmarks/HII\_region/} also holds a
\code{lexington\_cloudy.txt} comparison profile.  Nothing here was re-run.

\paragraph{Net comparison with \MoCHII{}.}  On channels \Torus{} is the closest
of the four other codes to \MoCHII{} in what it solves and the farthest in what
it omits.  It has the density-dependent level populations inside the thermal
balance, as \MoCHII{} now does through the $(T, \nel)$ table, on transition
energies taken from observed levels as \MoCHII{}'s rebuilt fits are, and from
the same routine that makes its line output; it has the gas--grain collisional
term as an option; and it puts
grain thermal emission into the gas balance and returns the metal line photons
to the transport, neither of which \Mocassin{} or \MoCHII{} does.  Against that
it omits five channels \MoCHII{} carries ---
collisional-ionization cooling, \HeIII{} recombination cooling, $Z = 2$
free-free, metal free-free and \HeII{} photoheating --- caps every ion at ten
levels, and cannot be extended past C, N, O, Ne and S without editing Fortran,
since its atomic data are \code{select case} blocks rather than files.  On data
its collision strengths are a mixture whose $\OIII{}$ agrees with CHIANTI~v11
to $3.5\%$ and whose $\NI{}$ is an order of magnitude above \Mocassin{}'s, so
the systematic $15$--$22\%$ excess found above in \Mocassin{}'s largest metal
coolants does not carry over to \Torus{} unexamined.

\subsection{Comparison, and one omission common to all five codes}

\begin{table}[htbp]
\centering
\caption{Cooling terms of the five Monte Carlo codes.  \MoCHII{}'s H/He
inventory is \code{src/cooling\_mod.f90}:225--258.  Ion counts: Wood, \CMI{}
and \Torus{} are fixed lists compiled into the source; the \Mocassin{} entry is
what its data files supply for
the HII40 composition (C, N, O, Ne, S) at the default \code{nstages = 7},
reaching O\,\textsc{vi}; \MoCHII{}'s 28 are its cooling fits, \HI{} plus 27
metal ions over 11 elements, reaching stage \textsc{iv} (\textsc{v} for Ca).  \MoCHII{} is a superset of the line and
H/He channels except for the last two rows: the gas--grain term,
which \Mocassin{} carries and \Torus{} carries as an option; and \HeI{}
collisional excitation, which no code carries.  It now carries the
density-dependent level populations in the thermal balance too, through the
$(T, \nel)$ table of Section~\ref{sec:limit}.}
\label{tab:codes}
\footnotesize
\begin{tabular}{lccccc}
\toprule
Cooling term & Wood & \CMI{} & \Mocassin{} & \Torus{} & \MoCHII{} \\
\midrule
metal collisionally excited lines & 12 ions & 13 ions & 26 ions & 15 ions & 28 ions \\
Cl, Ar line cooling      & data unread & no & yes & no & yes \\
\NI{}, \OI{} line cooling & defective & yes & yes & yes & yes \\
\SIV{} line cooling      & no  & yes & yes & yes & yes (added here) \\
\HI{} collisional excitation & no & no & yes & yes & yes \\
collisional-ionization cooling & no & no & no & no & yes \\
\HeIII{} recombination cooling & no & no & yes & no & yes \\
free-free at $Z = 2$ (\HeIII{}) & no & no & yes & no & yes \\
free-free from the metal ions & no & no & no & no & yes \\
level populations at the local $\nel$ & yes & yes & yes & yes & yes (table) \\
gas--grain collisional cooling & no & no & yes & optional & no \\
\HeI{} collisional excitation & no & no & no & no & no \\
\bottomrule
\end{tabular}
\end{table}

Table~\ref{tab:codes} collects the result.  The negative result of the audit is
the last row: \HeI{} collisional-excitation line cooling is absent from all
five codes, \MoCHII{} included.  In the HII40 budget of
Table~\ref{tab:budget} \Cloudy{} puts $1.943\times10^{35}$~erg\,s$^{-1}$ in it,
$0.062\%$ of the total, so at benchmark conditions the omission is harmless; it
would matter only where neutral helium is abundant and the gas hot enough to
excite it.

Ranking the Wood findings by what they are worth at HII40, using the \Cloudy{}
channel shares of Table~\ref{tab:budget}: the $\NIII{}$ collision-strength
inconsistency is the largest, a factor $2.1$ on $1.89\%$ of the budget; the
missing \HI{} collisional excitation is $0.72\%$; the unread chlorine and argon
blocks are $0.093\%$ and $0.006\%$; the missing \HeIII{} recombination cooling
is $0.001\%$ at this excitation.  The \NI{}/\OI{} defect is negligible at HII40
($[\OI{}]$ is $0.10\%$, $[\NI{}]$ $0.014\%$) but it lives precisely in the
partly neutral gas, so it matters for the ionization-front edge, for diffuse
ionized gas, and for anything PDR-like.  The $\OIII{}$ $^1S_0$ energy slip does
not move the cooling but biases $[\OIII{}]$ 4363 and the temperature diagnosed
from it by about $2\%$.

%==========================================================================
\section{Summary}
\label{sec:summary}

The \MoCHII{} temperature excess on the Lexington/Meudon benchmarks, which had
survived the elimination of charge exchange, reference convergence, spatial
resolution, recombination coefficients, collision-strength vintage, band
discretization, transport treatment and trace-species coupling, was two
defects in the offline construction of the Tier-1 cooling fits: ground-term
level populations set by Boltzmann factors, which at $kT \gg \Delta E_{\rm fs}$
is the high-density limit and suppressed the far-infrared fine-structure
cooling by $23.9\%$ of the total, and the Burgess--Tully scaling energy used
as the physical transition energy in both the excitation rate and the radiated
photon, worth a further $10.6\%$.  The diagnosis needed a fixed-state
comparison, because in a converged model the temperature rises until the books
balance and hides the deficit; the decisive internal evidence was that
\MoCHII{}'s own $n$-level solve and its own cooling fits disagreed by $28\%$
on an identical state.  Rebuilding all 28 fits from the $\nel \to 0$ limit of
statistical equilibrium with observed level energies, and adding $\SIV{}$,
moved HII40 inside the published temperature range and closed the photon
budget, the outer radius and the line ratios along with it.  That zero-density
limit still overcooled the low-critical-density lines at $\nH = 100$;
suppressing $\Lambda$ to the local $\nel$ with a $(T, \nel)$ table built from
the same $n$-level solve --- now the default,
\code{par\%cooling\_model = local\_ne} --- removes the far-infrared overestimate
and raised the converged benchmarks at once to 8211 and 7025~K, matching
\Cloudy{} c25.00 (8210, 6910~K); the zero-density fits are kept as the
\code{low\_density} option.  Later corrections to the stellar and diffuse
radiation field, none of them in the cooling, brought the production figures to
their current values: $8169$~K at HII40, inside the published
$7720$--$8199$~K range and $0.50\%$ below \Cloudy{}; and $6925$~K at HII20,
above the published $6402$--$6749$~K range but within $0.22\%$ of \Cloudy{}
c25.00's $6910$~K, which is itself $2.4\%$ above that range.  The HII20 range is
a 1995-era nine-code spread that both current-generation codes exceed, so the
verdict there is agreement with \Cloudy{}, not a shortfall against the
published floor.

The same audit applied to the other Monte Carlo photoionization codes shows
that \MoCHII{}'s inventory is now a superset of all four in the line and H/He
channels, with two exceptions.  \HeI{} collisional-excitation line cooling is
absent from Wood's code, from \CMI{}, from \Mocassin{}, from \Torus{} and from
\MoCHII{}, and is worth $0.062\%$ of the HII40 budget.  And \Mocassin{} carries a
gas--grain collisional term and \Torus{} carries one as an option, of order
$10^{-4}$ of the line cooling at
$\nH = 100$~cm$^{-3}$.  The metal level populations at the local $\nel$ that
the other four codes solve inside the thermal balance \MoCHII{} now matches
through the $(T, \nel)$ table (Section~\ref{sec:limit}), so that is no longer a
difference.  It also turned up defects on their side --- Wood's \NI{} and
\OI{} cooling is zeroed by an overwrite of the collision strengths, its cooling
and diagnostic routines disagree by a factor $2.1$ on the $\NIII{}$
57.3~$\mu$m collision strength (CHIANTI v11 supports the diagnostic value), its
chlorine and argon data are present but never read, its \SIV{} output slot is
never assigned, and its $\OIII{}$ $^1S_0$ energy is low by 113.7~cm$^{-1}$.
\CMI{} carries none of those, and the $11\%$ He$^+$ difference between the two
is in the recombination \emph{cooling} coefficient, not the recombination rate,
which they share exactly.  \Torus{} carries a Fortran exponent typo, \code{1.03-6}
for $1.03\times10^{-6}$, that makes one $\SII{}$ $A$-value negative and zeroes
three $[\SII{}]$ lines while inflating three others by $3.2$, and it enters the
$\NII{}$ $^1S_0$ statistical weight as $5$ instead of $1$.

\paragraph{Where the numbers come from.}  The \Cloudy{} reference runs and
their output files are in \code{tests/cloudy\_c25/} (\code{hii40\_cool.in} for
the cooling inventory, \code{hii\{40,20\}\_c25.in} for the line lists and
structure); the benchmark runs after the fix are
\code{tests/g2\_hii/hii\{40,20\}\_bench.in} with their \code{\_rates.h5} and
\code{\_lines.txt} output; the coefficient files are
\code{data/atomic/cooling\_<ion>.txt} and \code{nlevel\_<ion>.txt}, each with
its own provenance header; the fitting pipeline is
\code{tools/fitting/\{chianti\_cooling,fit\_cooling,fit\_nlevel\}.py}; and the
benchmark table columns are assembled by
\code{paper/tables/make\_wood\_tables.py}.  The channel-by-channel budget, the
defect isolation and the temperature closure were computed by analysis scripts
kept in the working session rather than committed, since each of them is a
one-off reading of the files listed above.

\end{document}
